\documentclass[11pt]{book}

% ==== Packages ====
\usepackage[utf8]{inputenc}
\usepackage[T1]{fontenc}
\usepackage{amsmath, amssymb, amsthm}
\usepackage{graphicx}
\usepackage{booktabs}
\usepackage{multirow}
\usepackage{array}
\usepackage{float}
\usepackage{hyperref}
\usepackage{cleveref}
\usepackage{xcolor}
\usepackage{tikz}
\usepackage{pgfplots}
\pgfplotsset{compat=1.18}
\usepackage{algorithm}
\usepackage{algorithmic}
\usepackage{geometry}
\usepackage{fancyhdr}
\usepackage{tocloft}
\usepackage{enumitem}
\usepackage{listings}
\usepackage{url}
\usepackage{cite}
\usepackage{appendix}
\usepackage{makeidx}
\makeindex

% ==== Page Setup ====
\geometry{a4paper, margin=1in}
\setlength{\headheight}{26pt}
\pagestyle{fancy}
\fancyhf{}
\fancyhead[LE,RO]{\thepage}
\fancyhead[LO]{\rightmark}
\fancyhead[RE]{\leftmark}

% ==== Colors ====
\definecolor{myblue}{RGB}{13,71,161}
\definecolor{mygreen}{RGB}{27,94,32}
\definecolor{myorange}{RGB}{230,81,0}
\definecolor{myviolet}{RGB}{94,53,177}
\definecolor{mygray}{RGB}{97,97,97}

% ==== Math Commands ====
\newcommand{\grad}{\nabla}
\newcommand{\EE}{\mathbb{E}}
\newcommand{\dd}{\mathrm{d}}
\newcommand{\kbT}{k_B T}
\newcommand{\vx}{\mathbf{x}}
\newcommand{\vz}{\mathbf{z}}
\newcommand{\vtheta}{\boldsymbol{\theta}}
\newcommand{\vphi}{\boldsymbol{\phi}}
\newcommand{\veps}{\boldsymbol{\epsilon}}
\newcommand{\KL}{\text{KL}}
\newcommand{\ELBO}{\text{ELBO}}
\newcommand{\JS}{\text{JS}}
\newcommand{\W}{\mathcal{W}}
\newcommand{\D}{\mathcal{D}}
\newcommand{\G}{\mathcal{G}}

% Bold vector/time-indexed variables
\newcommand{\vxz}{\mathbf{x}_0}
\newcommand{\vxt}{\mathbf{x}_t}
\newcommand{\alpt}{\alpha_t}
\newcommand{\sigt}{\sigma_t}
\newcommand{\vF}{\mathbf{F}}
\newcommand{\vinst}{\mathbf{v}}
\newcommand{\uavg}{\overline{\mathbf{u}}}
\newcommand{\dv}[2]{\frac{d#1}{d#2}}
\newcommand{\td}{\frac{d}{dt}}
\newcommand{\x}{\mathbf{x}}
\newcommand{\coloneqq}{\mathrel{\vcenter{\baselineskip0.5ex\lineskiplimit0pt\hbox{\scriptsize.}\hbox{\scriptsize.}}}=}
\newcommand{\sg}{\mathrm{sg}}
\newcommand{\JVP}{\mathrm{JVP}}

% ==== Theorem Environments ====
\theoremstyle{definition}
\newtheorem{definition}{Definition}[chapter]
\newtheorem{theorem}{Theorem}[chapter]
\newtheorem{lemma}{Lemma}[chapter]
\newtheorem{corollary}{Corollary}[chapter]
\newtheorem{proposition}{Proposition}[chapter]
\newtheorem{example}{Example}[chapter]
\newtheorem{remark}{Remark}[chapter]

% ==== Hyperref Setup ====
\hypersetup{
    colorlinks=true,
    linkcolor=myblue,
    urlcolor=myblue,
    citecolor=mygreen,
    bookmarksnumbered=true,
    bookmarksopen=true
}

% ==== Title Page ====
\title{Large Foundation Models: Mathematics, Algorithms, and Applications}
\author{Pipi Hu}
\date{\today}

\begin{document}

% ==== Front Matter ====
\frontmatter

% Title page
\maketitle

% Table of contents
\tableofcontents

% List of figures
\listoffigures

% List of tables
\listoftables

% Preface
\chapter*{Preface}
\addcontentsline{toc}{chapter}{Preface}

This book presents a comprehensive mathematical treatment of large foundation models, covering the theoretical foundations, algorithmic developments, and practical applications that have revolutionized artificial intelligence in recent years. The material is based on a course taught at the Beijing Institute of Mathematical Sciences and Applications (BIMSA).

The book is organized into three main parts:
\begin{itemize}
    \item \textbf{Part I: Generative Models} covers the mathematical foundations of modern generative modeling, including diffusion models, flow matching, variational autoencoders, normalizing flows, and generative adversarial networks.
    \item \textbf{Part II: Autoregressive Models} focuses on sequential generation methods, particularly transformer-based architectures and their applications in language modeling.
    \item \textbf{Part III: Advanced Topics} explores world models, reinforcement learning, reasoning, agents, and applications in AI for science.
\end{itemize}

Each chapter provides rigorous mathematical derivations, intuitive explanations, and practical insights that will be valuable for researchers, students, and practitioners working in the field of artificial intelligence.

% ==== Main Content ====
\mainmatter

% Part I: Generative Models
\part{Generative Models}

\chapter{Diffusion Models: From Mathematical Foundations to Practical Implementation}

\section{Introduction}

Diffusion models have emerged as one of the most powerful and versatile approaches to generative modeling in recent years. These models, inspired by non-equilibrium thermodynamics and statistical physics, have demonstrated remarkable success in generating high-quality images, audio, and other complex data types. The fundamental idea behind diffusion models is elegant yet profound: instead of directly modeling complex data distributions, we gradually transform simple noise into the target data through a carefully designed diffusion process that can be learned and reversed.

The mathematical foundations of diffusion models rest on several key concepts from probability theory, stochastic differential equations, and statistical physics. The score function, which represents the gradient of the log-probability density, plays a central role in these models. By learning to estimate this score function, we can generate samples from complex distributions using techniques like Langevin dynamics and other sampling methods.

This chapter provides a comprehensive mathematical treatment of diffusion models, starting from the fundamental principles of sampling from probability distributions and progressing through the theoretical foundations to practical implementations. We will explore the connections between different formulations, including score matching, denoising diffusion probabilistic models (DDPM), and continuous-time formulations. The chapter concludes with an examination of recent advances that have made diffusion models more efficient and practical.

\section{Preliminary: Generating Samples from Probability Distributions}

\subsection{The Challenge of Sampling}

Given a probability distribution $p(x)$, the fundamental question in generative modeling is how to sample from this distribution. For simple distributions like uniform or Gaussian distributions, sampling is straightforward using standard random number generators. For mixture Gaussian distributions, we can first determine which component to sample from and then sample from that specific Gaussian component.

However, for general probability distributions, especially those that are only known up to a normalization constant, sampling becomes significantly more challenging. This is where sophisticated sampling methods become essential. The most commonly used approaches include Langevin Markov Chain Monte Carlo (Langevin MCMC) sampling and Stein Variational Gradient Descent (SVGD), among others.

\subsection{The Role of the Score Function}

A typical process for generating new samples from a modeled data distribution involves several key steps. First, we need to find or approximate the probability function $p(x)$ by modeling $p(x;\theta)$ with parameters $\theta$. Second, we often work with unnormalized distributions of the form $p(x;\theta) = \frac{1}{C(\theta)}e^{-E(x;\theta)}$ where $C(\theta) \equiv \int_x e^{-E(x;\theta)}dx$ is the normalization constant. Finally, we generate samples using appropriate sampling methods.

The Langevin Markov Chain Monte Carlo sampling method provides a particularly elegant solution to this problem. The discrete-time Langevin dynamics can be written as:
\begin{equation}
x_t = x_{t-1} + \frac{\Delta t}{2}\nabla_x\log p(x_{t-1};\theta) + \sqrt{\Delta t}\epsilon, \quad \epsilon\sim N(0, 1).
\end{equation}

The crucial insight here is that the gradient term $\nabla_x\log p(x_{t-1};\theta) = -\nabla_x E(x;\theta)$ is independent of the normalization constant $C(\theta)$. This makes the method particularly attractive for sampling from unnormalized distributions.

We define the score function as:
\begin{equation}
S(x;\theta) \equiv \nabla \log p(x;\theta).
\end{equation}

The score function represents the direction of steepest increase in the log-probability density and plays a central role in many generative modeling approaches.

\subsection{Continuous-Time Formulation}

In the continuous-time limit as $\Delta t \to 0$, the Langevin MCMC process becomes:
\begin{equation}
dx = \frac{1}{2}\nabla_x\log p(x) dt + dB_t,
\end{equation}
where $dB_t$ is a Brownian motion. Given a distribution of the form $p(x) = \frac{1}{C}e^{-E}$, we have:
\begin{equation}
dx = -\frac{1}{2}\nabla_x E dt + dB_t.
\end{equation}

This continuous formulation reveals the connection to Langevin dynamics in molecular dynamics simulations. The full Langevin dynamics equation is:
\begin{equation}
\ddot{x} = -\nabla_x E - \gamma \dot{x} + dB_t,
\end{equation}
where $\gamma$ is a friction coefficient. When the friction term $\gamma \dot{x}$ dominates the acceleration term $\ddot{x}$, we obtain the overdamped Langevin equation:
\begin{equation}
dx = -\frac{1}{\gamma}\nabla_x E + \frac{1}{\gamma}dB_t.
\end{equation}

This overdamped limit is exactly what we use in the context of diffusion models and score-based generative modeling.

\section{Matching Score for Training}

\subsection{The Challenge of Score Matching}

The key insight in score-based generative modeling is to learn the score function using neural networks. The explicit score matching (ESM) loss is defined as:
\begin{equation}
J_{ESM} = \frac{1}{2} \mathbb{E}_p\left[\|s(x;\theta) - \frac{\partial \log p(x)}{\partial x}\|^2\right],
\end{equation}
where $s(x;\theta)$ is our neural network approximation of the true score function.

However, this formulation is not tractable in practice because we cannot directly compute $\frac{\partial \log p(x)}{\partial x}$ from data. This is where the theoretical developments in score matching become crucial.

\subsection{Implicit Score Matching}

Implicit Score Matching (ISM) was developed to make score matching tractable. The ISM loss is defined as:
\begin{equation}
J_{ISM} \equiv \mathbb{E}_p\left[\frac{1}{2}\|s(x;\theta)\|^2 + \nabla\cdot s(x;\theta)\right] = J_{ESM} - C,
\end{equation}
where $C$ is a constant independent of the parameters $\theta$.

The key insight in the proof of this equivalence involves expanding the ESM loss:
\begin{equation}
J_{ESM} = \frac{1}{2} \mathbb{E}_p\left[\|s(x;\theta) - \frac{\partial \log p(x)}{\partial x}\|^2\right]
\end{equation}
\begin{equation}
= \frac{1}{2} \mathbb{E}_p\left[\|s\|^2\right] - \mathbb{E}_p\left[s \cdot \frac{\partial \log p}{\partial x}\right] + \frac{1}{2}\mathbb{E}_p\left[\left\|\frac{\partial \log p}{\partial x}\right\|^2\right].
\end{equation}

The crucial step is to simplify the cross term using integration by parts:
\begin{equation}
- \mathbb{E}_p\left[s \cdot \frac{\partial \log p}{\partial x}\right] = -\int_x ps \cdot \nabla \log p \, dx
\end{equation}
\begin{equation}
= -\int_x ps \cdot \frac{\nabla p}{p} dx = -\int_x s \cdot \nabla p \, dx
\end{equation}
\begin{equation}
= -\int \nabla \cdot (ps) dx + \int p \nabla \cdot s \, dx = \mathbb{E}_p[\nabla \cdot s].
\end{equation}

Therefore, we have $J_{ISM} = J_{ESM} - C$ where $C \equiv \frac{1}{2}\mathbb{E}_p\left[\left\|\frac{\partial \log p}{\partial x}\right\|^2\right]$ is a constant independent of $\theta$. This makes ISM tractable since we don't need to compute the true score function.

\subsection{Denoising Score Matching}

However, the ISM loss is not very stable to optimize for two main reasons. First, the expectation is performed over the entire distribution, which can be computationally expensive. Second, the loss can become negative and decrease to a large negative value like $-10^5$, making optimization challenging.

Denoising Score Matching (DSM) was developed to address these problems. The DSM loss is defined as:
\begin{equation}
J_{DSM} = \frac{1}{2}\mathbb{E}_{p(x,\tilde{x})}\left[\|s(\tilde{x};\theta) - \nabla_{\tilde{x}}\log p(\tilde{x}|x)\|^2\right].
\end{equation}

We can prove that $J_{DSM} = J_{ESM} + C$, where $C$ is a constant. The key insight is to show that the cross terms of $J_{ESM}$ and $J_{DSM}$ are equal:
\begin{equation}
\mathbb{E}_{p(\tilde{x})}[s(\tilde{x}) \cdot \nabla_{\tilde{x}}\log p(\tilde{x})] = \mathbb{E}_{p(\tilde{x},x)}\left[s(\tilde{x}) \cdot \nabla_{\tilde{x}} \log p(\tilde{x}|x)\right].
\end{equation}

This equivalence is established by expanding the left side and using the relationship between the marginal and conditional distributions.

However, a problem arises when $\tilde{x} \to x$. We can prove that when Gaussian noise is added by $p(\tilde{x}|x) = N(\alpha x, \sigma^2)$, the DSM loss diverges to infinity as $\tilde{x} \to x$. This is because:
\begin{equation}
\nabla_{\tilde{x}}\log p(\tilde{x}|x) = -\frac{\tilde{x}-\alpha x}{\sigma^2},
\end{equation}
and the expectation $\mathbb{E}_{p(x,\tilde{x})}[\|\nabla_{\tilde{x}}\log p(\tilde{x}|x)\|^2]$ goes to infinity as $\sigma \to 0$.

\section{Relations Between Different Models}

\subsection{Gaussian Noise Addition}

Let us revisit the process of adding noise from $x$ to $\tilde{x}$. The Gaussian transition kernel for adding noise is:
\begin{equation}
\tilde{x} = \alpha x + \sigma \epsilon,
\end{equation}
which is equivalent to the Gaussian transition kernel:
\begin{equation}
p(\tilde{x}|x) = N(\tilde{x}; \alpha x, \sigma^2),
\end{equation}
where $\epsilon \sim N(0,1)$ and:
\begin{equation}
N(\tilde{x};\alpha x, \sigma^2) = C e^{-\frac{\|\tilde{x}-\alpha x\|^2}{2\sigma^2}}.
\end{equation}

We also have:
\begin{equation}
\epsilon = \frac{\tilde{x}-\alpha x}{\sigma}.
\end{equation}

The question naturally arises: how can we recover $x$ given $\tilde{x}$?

\subsection{Tweedie's Formula}

Tweedie's Formula from Bayesian statistics provides a crucial connection. Given random variables $x, y$ where $y \sim N(x, \sigma^2 I)$, i.e., $p(y|x) = N(x, \sigma^2)$, the expectation of $x$ can be given by:
\begin{equation}
\mathbb{E}[x|y] = y + \sigma^2 \nabla \log p(y).
\end{equation}

Letting $x \leftarrow \alpha x$ and $y \leftarrow \tilde{x}$, we have:
\begin{equation}
s \equiv \nabla_{\tilde{x}} \log p(\tilde{x}) = -\frac{\tilde{x}-\alpha \mathbb{E}[x|\tilde{x}]}{\sigma^2} = -\frac{1}{\sigma}\mathbb{E}\left[\frac{\tilde{x}-\alpha x}{\sigma}|\tilde{x}\right] = -\frac{\mathbb{E}[\epsilon|\tilde{x}]}{\sigma},
\end{equation}
where we have used the fact that $\epsilon = \frac{\tilde{x}-\alpha x}{\sigma}$.

\subsection{Alternative Derivation from Score Definition}

We can also derive this relationship directly from the definition of the score function. Given the noise addition process $\tilde{x} = \alpha x + \sigma \epsilon$ and the conditional distribution:
\begin{equation}
p(\tilde{x}|x) = Ce^{-\frac{\|\tilde{x}-\alpha x\|^2}{2\sigma^2}},
\end{equation}
we have:
\begin{equation}
s(\tilde{x}) = \nabla_{\tilde{x}}\log p(\tilde{x}) = \frac{\nabla_{\tilde{x}}p(\tilde{x})}{p(\tilde{x})} = \frac{\int_x \nabla_{\tilde{x}}p(\tilde{x}|x) p(x)dx}{p(\tilde{x})}.
\end{equation}

From the conditional distribution, we have:
\begin{equation}
s(\tilde{x}) = \int_x \nabla_{\tilde{x}} \log p(\tilde{x}|x)p(x|\tilde{x})dx = \mathbb{E}_x[\nabla_{\tilde{x}} \log p(\tilde{x}|x)| \tilde{x}] = \mathbb{E}_x\left[-\frac{\tilde{x}-\alpha x}{\sigma^2}| \tilde{x}\right] = -\frac{\tilde{x}-\alpha \mathbb{E}[x|\tilde{x}]}{\sigma^2}.
\end{equation}

\subsection{Theoretical Foundation: Conditional Expectation}

The theoretical foundation for these relationships lies in the theory of conditional expectation. Let $X$ be an integrable random variable. Then for each $\sigma$-algebra $\mathcal{V}$ and $Z \in \mathcal{V}$, $Z^* = \mathbb{E}(X| \mathcal{V})$ solves the least square problem:
\begin{equation}
\text{argmin}_{Z\in \mathcal{V}} \| Z - X\|,
\end{equation}
where $\| Z\| = \left( \int Z^2 dP \right)^{\frac{1}{2}}$.

This result establishes that the conditional expectation is the optimal predictor in the mean-square sense, which provides the theoretical justification for the relationships we have derived.

\subsection{Three Equivalent Model Formulations}

Based on the above analysis, we can identify three equivalent ways to model the denoising process:

\textbf{Score model:} The loss function is:
\begin{equation}
\text{Loss}_{score} = \|s_\theta(\tilde{x}) - \frac{\epsilon}{\sigma}\|^2.
\end{equation}

\textbf{X prediction model (denoising model, $x_0$ model):} The loss function is:
\begin{equation}
\text{Loss}_{x_0} = \|x_\theta (\tilde{x}) - x\|^2,
\end{equation}
with the relationship:
\begin{equation}
s(\tilde{x}) = -\frac{\tilde{x}-\alpha x_{\theta^*} (\tilde{x})}{\sigma^2}
\end{equation}
where $x_\theta(\tilde{x}) \to x_{\theta^*}(\tilde{x}) \equiv \mathbb{E}[x|\tilde{x}]$.

\textbf{$\epsilon$ prediction model:} The loss function is:
\begin{equation}
\text{Loss}_{\epsilon} = \|\epsilon_\theta(\tilde{x}) - \epsilon\|^2,
\end{equation}
with the relationship:
\begin{equation}
s(\tilde{x}) = -\frac{\epsilon_{\theta^*}(\tilde{x})}{\sigma},
\end{equation}
where $\epsilon_\theta(\tilde{x}) \to \epsilon_{\theta^*}(\tilde{x}) \equiv \mathbb{E}[\epsilon|\tilde{x}]$.

In the optimal case ($\theta^*$), these three formulations are consistent:
\begin{equation}
s_{\theta^{*}}(\tilde{x}) = -\frac{\tilde{x}-\alpha x_{\theta^*} (\tilde{x})}{\sigma^2} = -\frac{\epsilon_{\theta^*}(\tilde{x})}{\sigma}.
\end{equation}

In practice, we use the connection between these three models without necessarily reaching the optimal parameters:
\begin{equation}
s_\theta(\tilde{x}) = -\frac{\tilde{x}-\alpha x_{\theta} (\tilde{x})}{\sigma^2} = -\frac{\epsilon_{\theta}(\tilde{x})}{\sigma},
\end{equation}
and by substituting each form into the loss definitions, we can recover all the losses given by different models.

\section{Potential Score Matching Loss}

\subsection{Key Identity: Jacobian Switch}

The Gaussian transition density is:
\begin{equation}
q(\mathbf{x}_t \mid \mathbf{x}_0) = \mathcal{N}(\mathbf{x}_t; \alpha_t \mathbf{x}_0, \sigma_t^2\mathbf{I}) \propto \exp\left(-\frac{\|\mathbf{x}_t-\alpha_t \mathbf{x}_0\|^2}{2\sigma_t^2}\right).
\end{equation}

A basic but crucial identity is:
\begin{equation}
\nabla_{\mathbf{x}_t} q(\mathbf{x}_t \mid \mathbf{x}_0) = -\frac{1}{\alpha_t} \nabla_{\mathbf{x}_0} q(\mathbf{x}_t \mid \mathbf{x}_0), \quad (\alpha_t > 0).
\end{equation}

The proof of this identity follows from differentiating the exponent with respect to $\mathbf{x}_t$ and $\mathbf{x}_0$:
\begin{equation}
\nabla_{\mathbf{x}_t}\|\mathbf{x}_t-\alpha_t\mathbf{x}_0\|^2 = 2(\mathbf{x}_t-\alpha_t\mathbf{x}_0), \quad \nabla_{\mathbf{x}_0}\|\mathbf{x}_t-\alpha_t\mathbf{x}_0\|^2 = -2\alpha_t(\mathbf{x}_t-\alpha_t\mathbf{x}_0),
\end{equation}
thus giving the factor $-1/\alpha_t$.

\subsection{Score as Conditional Force Expectation}

With $p(\mathbf{x}_0) \propto e^{-E(\mathbf{x}_0)/k_B T}$ and $\mathbf{x}_t = \alpha_t \mathbf{x}_0 + \sigma_t \boldsymbol{\epsilon}$, we have:
\begin{equation}
\nabla_{\mathbf{x}_t}\log p(\mathbf{x}_t) = \frac{1}{\alpha_t} \mathbb{E}_{\mathbf{x}_0 \mid \mathbf{x}_t}\left[\frac{\mathbf{F}(\mathbf{x}_0)}{k_B T}\right],
\end{equation}
where $\mathbf{F}(\mathbf{x}_0) = -\nabla_{\mathbf{x}_0}E(\mathbf{x}_0)$.

The derivation follows from:
\begin{align}
\nabla_{\mathbf{x}_t}\log p(\mathbf{x}_t) &= \frac{\nabla_{\mathbf{x}_t}\int p(\mathbf{x}_0) q(\mathbf{x}_t \mid \mathbf{x}_0) d \mathbf{x}_0}{p(\mathbf{x}_t)}\\
&= \frac{1}{p(\mathbf{x}_t)}\int p(\mathbf{x}_0) \nabla_{\mathbf{x}_t} q(\mathbf{x}_t \mid \mathbf{x}_0) d \mathbf{x}_0
\end{align}

Using the Jacobian switch identity:
\begin{align}
&= \frac{1}{p(\mathbf{x}_t)}\int p(\mathbf{x}_0)\left(-\frac{1}{\alpha_t}\right)\nabla_{\mathbf{x}_0} q(\mathbf{x}_t \mid \mathbf{x}_0) d \mathbf{x}_0\\
&= -\frac{1}{\alpha_t} \frac{1}{p(\mathbf{x}_t)}\int q(\mathbf{x}_t \mid \mathbf{x}_0) \nabla_{\mathbf{x}_0} p(\mathbf{x}_0) d \mathbf{x}_0
\end{align}

Finally:
\begin{align}
&= -\frac{1}{\alpha_t}\int \frac{p(\mathbf{x}_0) q(\mathbf{x}_t \mid \mathbf{x}_0)}{p(\mathbf{x}_t)} \nabla_{\mathbf{x}_0}\log p(\mathbf{x}_0) d \mathbf{x}_0\\
&= \frac{1}{\alpha_t} \mathbb{E}_{\mathbf{x}_0 \mid \mathbf{x}_t}\left[\frac{-\nabla_{\mathbf{x}_0}E(\mathbf{x}_0)}{k_B T}\right] = \frac{1}{\alpha_t} \mathbb{E}_{\mathbf{x}_0 \mid \mathbf{x}_t}\left[\frac{\mathbf{F}(\mathbf{x}_0)}{k_B T}\right]
\end{align}

This result shows that the score function can be interpreted as a conditional expectation of the force field, providing a physical interpretation of the diffusion process.

\section{From Theory to Practice: Implementation Approaches}

Now that we have established the theoretical foundations of score matching and diffusion processes, let us explore how these concepts are implemented in practice through several key approaches: Score Matching Langevin Dynamics (SMLD), Denoising Diffusion Probabilistic Models (DDPM), and continuous formulations.

\subsection{Score Matching Langevin Dynamics (SMLD)}

SMLD adds noise in the following manner:
\begin{equation}
p(\tilde{x}_i| x) = N(\tilde{x}_i; x, \sigma_i^2) \equiv Ce^{-\frac{\|\tilde{x}_i-x\|^2}{2\sigma_i^2}},
\end{equation}
where a geometric sequence $\sigma_1 > \sigma_2 > \cdots > \sigma_T \approx 0$ is used to add noise with different levels.

In terms of random variables:
\begin{equation}
\tilde{x}_i = x + \sigma_i \epsilon
\end{equation}
for different $\sigma_i$ to get different noised random variables $\tilde{x}_i$.

The training objective $J_{DSM}$ is:
\begin{equation}
J_{DSM} = \frac{1}{2}\mathbb{E}_{\sigma_i}\mathbb{E}_{p(x)}\mathbb{E}_{p(\tilde{x}_i|x)}\left[\|s_\theta(\tilde{x}_i,\sigma_i) + \frac{\tilde{x}_i-x}{\sigma_i^2}\|^2\right].
\end{equation}

For sampling, annealed Langevin MCMC is used:
\begin{equation}
x_t = x_{t-1} + \frac{\Delta t}{2}S_\theta(x_{t-1}) + \sqrt{\Delta t}\epsilon, \quad \epsilon \sim N(0,1).
\end{equation}

\subsection{Denoising Diffusion Probabilistic Models (DDPM)}

DDPM was derived from the Evidence Lower Bound (ELBO), not directly from score matching. The main procedure involves:

\textbf{Adding noise:}
\begin{equation}
x_t = \sqrt{1-\beta_t}x_{t-1} + \sqrt{\beta_t}\epsilon;
\end{equation}

\textbf{Transition kernel:}
\begin{equation}
p(x_t|x_0) = N(x_t;\alpha_t x_0, \sigma_t^2),
\end{equation}
where $\alpha_t = \prod_{i=1}^{t} \sqrt{1-\beta_i}$ and $\sigma_t^2 = 1-\alpha_t^2$.

\textbf{Training Loss:}
\begin{equation}
\mathbb{E}_{t \sim U[0,1]}\mathbb{E}_{p(x)}\mathbb{E}_{p(x_t|x)}\left[\|\epsilon_\theta(x_t,t) - \epsilon\|^2\right],
\end{equation}
where we use the approximation $x_0 = \frac{x_t-\sigma_t\epsilon}{\alpha_t}$.

\textbf{Sampling:}
\begin{equation}
x_{t-1} = \frac{1}{\sqrt{1-\beta_t}}(x_t + \beta_t s_\theta(x_t, t)) + \sqrt{\beta_t}\epsilon_t, \quad t = T, T-1, \ldots, 1.
\end{equation}

\section{Continuous Formulations}

\subsection{Continuous Version of SMLD}

Recall that $x_t = x_0 + \sigma_t\epsilon$. We can prove that:
\begin{equation}
x_t = x_{t-1} + \sqrt{\sigma_t^2-\sigma_{t-1}^2}\epsilon.
\end{equation}

The proof follows by recurrence:
\begin{align}
x_t &= x_{t-1} + \sqrt{\sigma_t^2-\sigma_{t-1}^2}\epsilon\\
&= x_{t-2} + \sqrt{\sigma_{t-1}^2-\sigma_{t-2}^2}\epsilon + \sqrt{\sigma_t^2-\sigma_{t-1}^2}\epsilon\\
&= x_{t-2} + \sqrt{\sigma_t^2-\sigma_{t-2}^2}\epsilon\\
&= \ldots = x_0 + \sigma_t\epsilon, \quad \text{where } \sigma_0 \approx 0.
\end{align}

From this, we have:
\begin{equation}
\Delta x_t = \sqrt{\frac{\sigma_t^2-\sigma_{t-1}^2}{\Delta t}}\sqrt{\Delta t}\epsilon = \sqrt{\frac{\sigma_t^2-\sigma_{t-1}^2}{\Delta t}}\Delta B_t.
\end{equation}

Taking the limit as $\Delta t \to 0$:
\begin{equation}
dx = \sqrt{\frac{d\sigma_t^2}{dt}}dB_t.
\end{equation}

\subsection{Continuous Version of DDPM}

Recall that $x_t = \sqrt{1-\beta_t}x_{t-1} + \sqrt{\beta_t}\epsilon$. Similarly:
\begin{equation}
x_t \approx (1-\frac{1}{2}\beta_t)x_{t-1} + \sqrt{\beta_t}\epsilon,
\end{equation}
hence:
\begin{align}
\Delta x_t &= -\frac{1}{2}\tilde{\beta_t} x_{t-1}\Delta t + \sqrt{\tilde{\beta_t}}\sqrt{\Delta t}\epsilon, \quad \text{where } \tilde{\beta_t} = \frac{\beta_i}{\Delta t}\\
&= -\frac{1}{2}\tilde{\beta_t} x_{t-1}\Delta t + \sqrt{\tilde{\beta_t}}\Delta B_t.
\end{align}

Taking the limit as $\Delta t \to 0$:
\begin{equation}
dx = -\frac{1}{2}\tilde{\beta_t} x dt + \sqrt{\tilde{\beta_t}}dB_t.
\end{equation}

\subsection{Variable Elimination (VE) and Variable Preservation (VP)}

Summarizing the forms of continuous DDPM and SMLD, we can write the adding noise process as:
\begin{equation}
dx = f(x, t)dt + g(t) dB_t,
\end{equation}
where the corresponding distribution $p(x, t)$ satisfies the Fokker-Planck equation:
\begin{equation}
\frac{\partial p(x,t)}{\partial t} + \nabla \cdot (f(x,t) p(x,t)) - \frac{1}{2}g^2(t) \triangle p(x,t) = 0.
\end{equation}

The training DSM objective is:
\begin{equation}
J_{DSM} = \mathbb{E}_{t \sim U[0,1]}\mathbb{E}_{p(x_0)}\mathbb{E}_{p(x_t|x_0)}\left[\|s_\theta(x_t,t) - \nabla_{x_t}\log p(x_t|x_0)\|^2\right].
\end{equation}

\subsection{Transition Kernels for VE and VP}

For VE, the transition kernel is:
\begin{equation}
P(x(t)|x(0)) = N(x(t); x(0), \sigma^2(t)-\sigma^2(0)),
\end{equation}
where $\tilde{\alpha}_t \equiv 1$ and $\tilde{\sigma}^2_t \equiv \sigma^2(t)-\sigma^2(0) \approx \sigma^2(t)$.

For VP, the transition kernel is:
\begin{equation}
p(x(t)|x(0)) = N\left(x(t); x(0)e^{-\frac{1}{2}\int_0^t \tilde{\beta}(s)ds}, 1-e^{-\int_0^t \tilde{\beta}(s)ds}\right),
\end{equation}
where $\tilde{\alpha}_t \equiv e^{-\frac{1}{2}\int_0^t \tilde{\beta}(s)ds}$ and $\tilde{\sigma}^2_t \equiv 1-e^{-\int_0^t \tilde{\beta}(s)ds}$.

\section{Training and Inference}

\subsection{Training Recipe}

The training for VE follows this process (VP is similar):

1. Randomly choose one data point $x \sim p_{data}$;
2. Randomly sample $t \sim U[0, 1]$;
3. Randomly sample white noise $\epsilon \sim N(0, 1)$;
4. Add noise: $x_t = \tilde{\alpha}_t x + \tilde{\sigma}_t\epsilon$, where the mean value $\tilde{\alpha}_t x = x(0)$ and standard deviation $\tilde{\sigma}_t = \sqrt{\sigma^2(t)-\sigma^2(0)} \approx \sigma(t)$ in VE;
5. Compute the loss by summing the following norm over $x$, $t$ and $\epsilon$:
\begin{equation}
\|s_\theta(x_t, t) - \epsilon/\tilde{\sigma}_t\|.
\end{equation}

\subsection{Inference: Reverse Denoising Process}

The reverse denoising process is given by:
\begin{equation}
dx = (f(x,t) - g^2\nabla_x \log p(x,t))dt + g(t)dB_t,
\end{equation}
or:
\begin{equation}
dx = (f(x,t) - \frac{1}{2}g^2\nabla_x \log p(x,t))dt,
\end{equation}
or:
\begin{equation}
dx = (f(x,t) - \frac{3}{2}g^2\nabla_x \log p(x,t))dt + \sqrt{2}g(t)dB_t.
\end{equation}

All these forms are equivalent and obey the same Fokker-Planck equation for $p(x, t)$.

\subsection{Inference Recipe}

The inference for VE follows this process (VP is similar):

1. Randomly generate noise at time $t=1$ as $x(1) \sim N(0, \sigma^2_{max})$;
2. Integrate the reverse sampling equation:
\begin{equation}
dx = \left(f(x,t) - g^2(t) s_\theta(x,t)\right)dt + g(t)dB_t
\end{equation}
or:
\begin{equation}
dx = (f - \frac{1}{2}g^2 s_\theta)dt
\end{equation}
over the time span $[0, 1]$ to get $x(0)$;
3. Corrector process: at each internal value $x(t)$, we can relax the value $x(t)$ by a corrector, which takes the Langevin MCMC process.

\subsection{The Corrector: Langevin MCMC}

The Langevin diffusion process is:
\begin{equation}
x_i = x_{i-1} + \frac{\Delta t}{2}\nabla_x\log p(x_{i-1}) + \sqrt{\Delta t}\epsilon, \quad \epsilon \sim N(0, 1).
\end{equation}

We can prove that the distribution $\pi(x)$ of the obtained data points $\{x_i\}_{i=1}^N$ converges to $p(x)$. The continuous form is:
\begin{equation}
dx = \frac{1}{2}\nabla\log p(x) dt + dB_t.
\end{equation}

The corresponding Fokker-Planck equation is:
\begin{equation}
\frac{\partial\pi(x,t)}{\partial t} + \nabla \cdot \left(\frac{1}{2}\nabla\log p(x) \pi(x,t)\right) - \frac{1}{2}\triangle \pi(x,t) = 0.
\end{equation}

With sufficient time evolution, the distribution converges to a stable distribution where:
\begin{equation}
\frac{\partial \pi(x,t)}{\partial t} = 0, \quad t \to \infty.
\end{equation}

Hence:
\begin{equation}
\nabla \cdot (\pi\nabla\log p - \nabla\pi) = 0, \quad \forall x.
\end{equation}

Therefore:
\begin{equation}
\nabla\log p = \nabla\log \pi
\end{equation}
given that $\pi > 0$ almost everywhere. Hence:
\begin{equation}
\pi^*(x) = p(x)
\end{equation}
where $\pi^*$ is the stable distribution of $\pi(x, t)$.

\section{Recent Advances and Practical Considerations}

\subsection{DDPM: Denoising Diffusion Probabilistic Models}

DDPM represents a foundational approach to diffusion models. The goal is to learn a generative model $p_\theta(x_0)$ by reversing a noisy diffusion process. The model introduces latents $x_{1:T}$ with the same dimensionality as data $x_0$.

The reverse (sampling) chain is:
\begin{equation}
p_\theta(x_{0:T}) = p(x_T)\prod_{t=1}^{T} p_\theta(x_{t-1} \mid x_t), \quad p(x_T) = \mathcal{N}(0,I),
\end{equation}
\begin{equation}
p_\theta(x_{t-1} \mid x_t) = \mathcal{N}(\mu_\theta(x_t,t), \Sigma_\theta(x_t,t)).
\end{equation}

The forward (fixed) diffusion is:
\begin{equation}
q(x_{1:T} \mid x_0) = \prod_{t=1}^{T} q(x_t \mid x_{t-1}), \quad q(x_t \mid x_{t-1}) = \mathcal{N}(\sqrt{1-\beta_t} x_{t-1}, \beta_t I).
\end{equation}

With $\alpha_t := 1-\beta_t$ and $\bar{\alpha}_t := \prod_{s=1}^t \alpha_s$:
\begin{equation}
q(x_t \mid x_0) = \mathcal{N}(\sqrt{\bar{\alpha}_t} x_0, (1-\bar{\alpha}_t)I).
\end{equation}

The training objective optimizes the negative log-likelihood via a reduced-variance ELBO:
\begin{align}
\mathbb{E}[-\log p_\theta(x_0)] &\leq \underbrace{\mathbb{E}[\text{KL}(q(x_T|x_0) \| p(x_T))]}_{L_T} \nonumber \\
&\quad + \sum_{t=2}^{T} \underbrace{\mathbb{E}[\text{KL}(q(x_{t-1}|x_t,x_0) \| p_\theta(x_{t-1}|x_t))]}_{L_{t-1}} \nonumber \\
&\quad - \underbrace{\mathbb{E}[\log p_\theta(x_0|x_1)]}_{L_0}.
\end{align}

The $\epsilon$-prediction parameterization writes the forward noising as:
\begin{equation}
x_t = \sqrt{\bar{\alpha}_t} x_0 + \sqrt{1-\bar{\alpha}_t} \varepsilon, \quad \varepsilon \sim \mathcal{N}(0,I).
\end{equation}

The reverse mean is parameterized via a network $\varepsilon_\theta(x_t,t)$:
\begin{equation}
\mu_\theta(x_t,t) = \frac{1}{\sqrt{\alpha_t}} \left(x_t - \frac{1-\alpha_t}{\sqrt{1-\bar{\alpha}_t}} \varepsilon_\theta(x_t,t)\right).
\end{equation}

The simplified objective is:
\begin{equation}
\mathcal{L}_{\text{simple}}(\theta) = \mathbb{E}_{t \sim \mathcal{U}[1,T], x_0, \varepsilon} \|\varepsilon - \varepsilon_\theta(\sqrt{\bar{\alpha}_t} x_0 + \sqrt{1-\bar{\alpha}_t} \varepsilon, t)\|_2^2.
\end{equation}

\subsection{EDM: Elucidating the Design Space}

EDM separates design choices (sampling, schedules, preconditioning, training) to simplify and speed up diffusion models. The key outcomes include faster sampling (tens of NFEs) with strong FID and modular choices for ODE view, schedule $\sigma(t)$, scaling $s(t)$, time steps $\{t_i\}$, preconditioning, and loss weighting.

The probability flow ODE with schedule $\sigma(t)$ and scaling $s(t)$ is:
\begin{equation}
\frac{d\mathbf{x}}{dt} = \left(\frac{\dot{\sigma}}{\sigma} + \frac{\dot{s}}{s}\right)\mathbf{x} - \frac{\dot{\sigma} s}{\sigma} \nabla_{\mathbf{x}}\log p\left(\frac{\mathbf{x}}{s};\sigma\right).
\end{equation}

Using denoising score matching:
\begin{equation}
\nabla_{\mathbf{x}}\log p(\mathbf{x};\sigma) = \frac{D(\mathbf{x};\sigma) - \mathbf{x}}{\sigma^2},
\end{equation}
so the ODE is solved numerically by evaluating a denoiser $D_\theta(\mathbf{x};\sigma)$.

\subsection{DPM-Solver: Fast ODE Solver}

DPM-Solver exploits the semi-linear structure of the diffusion ODE to deliver high-quality samples in 10-20 NFE. The key idea is to compute the linear part analytically and only approximate a special, exponentially weighted integral of the network.

The forward SDE (VP-type example) is:
\begin{equation}
d\mathbf{x}_t = f(t)\mathbf{x}_t dt + g(t) d\mathbf{w}_t,
\end{equation}
where $f(t) = \frac{d}{dt}\log \alpha_t$ and $g^2(t) = \frac{d}{dt}\sigma_t^2 - 2\frac{d}{dt}\log \alpha_t \cdot \sigma_t^2$.

The probability flow ODE with learned noise predictor $\varepsilon_\theta(\mathbf{x},t)$ is:
\begin{equation}
\frac{d\mathbf{x}_t}{dt} = f(t)\mathbf{x}_t + \frac{g^2(t)}{2\sigma_t} \varepsilon_\theta(\mathbf{x}_t,t), \quad \mathbf{x}_T \sim \mathcal{N}(0,\tilde{\sigma}^2\mathbf{I}).
\end{equation}

The exact solution via variation of constants is:
\begin{equation}
\mathbf{x}_t = e^{\int_t^s f(\tau) d\tau} \mathbf{x}_s + \int_t^s e^{\int_t^\tau f(r) dr} \frac{g^2(\tau)}{2\sigma_\tau} \varepsilon_\theta(\mathbf{x}_\tau,\tau) d\tau.
\end{equation}

Introducing half-logSNR $\lambda_t := \log(\alpha_t/\sigma_t)$ and using change of variables:
\begin{equation}
\mathbf{x}_t = \frac{\alpha_t}{\alpha_s} \mathbf{x}_s - \alpha_t \int_{\lambda_s}^{\lambda_t} e^{-\lambda} \hat{\varepsilon}_\theta(\hat{\mathbf{x}}_\lambda,\lambda) d\lambda.
\end{equation}

\subsection{Consistency Models}

Consistency models enable fast 1-step generation while keeping diffusion model perks (quality, zero-shot editing, compute/quality trade-off). The idea is to learn a map that sends any point on a Probability-Flow (PF) ODE trajectory back to its origin (data).

The PF ODE from SDE is:
\begin{equation}
\frac{dx_t}{dt} = -t s_\phi(x_t,t),
\end{equation}
where $s_\phi \approx \nabla \log p_t$ under EDM-style settings.

A consistency function $f:(x_t,t) \mapsto x_\varepsilon$ with self-consistency $f(x_t,t) = f(x_{t'},t')$ if on the same trajectory. A consistency model is a neural estimator $f_\theta(x,t)$ of $f$ trained to enforce self-consistency.

For one-step generation:
\begin{equation}
x_T \sim \mathcal{N}(0,T^2 I), \quad x_\varepsilon = f_\theta(x_T,T).
\end{equation}

For few-step generation, we alternate between denoising and noise injection:
\begin{equation}
x \leftarrow f_\theta(x_T,T), \text{ then for } \tau_1 > \ldots > \tau_{N-1}: \hat{x}_{\tau_n} \leftarrow x + \sqrt{\tau_n^2-\varepsilon^2} z, z \sim \mathcal{N}(0,I); x \leftarrow f_\theta(\hat{x}_{\tau_n},\tau_n).
\end{equation}

\section{Open Questions and Future Directions}

Several important questions remain open in the field of diffusion models:

1. \textbf{Equivalence of reverse forms:} Please prove that the forms of the reverse process are equivalent with the same marginal distribution $p(x,t)$ given the same initial condition.

2. \textbf{Score function dynamics:} Please prove that the score function $s(x,t)$ satisfies:
\begin{equation}
\frac{\partial s}{\partial t} = \nabla \left(-\nabla \cdot f - \langle f, s\rangle + \frac{1}{2}g^2\|s\|^2 + \frac{1}{2}g^2\langle \nabla, s \rangle\right).
\end{equation}

3. \textbf{Practical considerations:} Why is ISM loss not commonly used like DSM loss in diffusion models nowadays? This question invites discussion about the practical trade-offs between different loss formulations.

These questions represent active areas of research that could lead to further theoretical understanding and practical improvements in diffusion models.

\section{Conclusion}

Diffusion models represent a powerful and theoretically well-founded approach to generative modeling. The mathematical foundations, rooted in stochastic differential equations, score matching, and statistical physics, provide a solid theoretical basis for understanding these models. The connections between different formulations—score matching, DDPM, and continuous-time approaches—reveal the underlying unity of these methods.

The practical implementations, from SMLD to DDPM to modern approaches like EDM and consistency models, demonstrate the evolution of the field toward more efficient and practical methods. The theoretical understanding of the relationships between different model formulations provides valuable insights for practitioners and researchers.

As the field continues to evolve, the mathematical foundations established in this chapter will remain relevant for understanding new developments and for designing improved methods. The open questions identified provide directions for future research that could lead to further theoretical insights and practical improvements.

The success of diffusion models in generating high-quality samples across diverse domains---from images to audio to scientific data---demonstrates the power of this approach. The mathematical rigor underlying these methods ensures that they will continue to be a valuable tool in the arsenal of generative modeling techniques.

\chapter{Flow Matching: A Unifying Perspective on Generative Modeling}

\section{Introduction}

Flow matching represents a powerful and mathematically elegant approach to generative modeling that unifies various perspectives on learning to transform simple distributions into complex data distributions. Unlike score-based methods that require estimating gradients of log-probability densities, flow matching provides direct supervision through analytically derived velocity fields.

The fundamental idea behind flow matching is to learn a velocity field $\mathbf{v}(\mathbf{x},t)$ that transports a simple base distribution $p_1(\mathbf{x})$ to the data distribution $p_0(\mathbf{x})$ through the continuity equation:
\begin{equation}
    \partial_t p(\mathbf{x},t) + \nabla \cdot (p(\mathbf{x},t) \mathbf{v}(\mathbf{x},t)) = 0.
\end{equation}

This chapter provides a comprehensive mathematical treatment of flow matching, starting from the theoretical foundations and progressing through practical implementations and recent advances.

\section{Mathematical Foundations}

\subsection{The Continuity Equation}

The continuity equation is the fundamental mathematical principle underlying flow matching. Given a velocity field $\mathbf{v}(\mathbf{x},t)$, the probability density $p(\mathbf{x},t)$ transported by the ordinary differential equation $\dot{\mathbf{x}} = \mathbf{v}(\mathbf{x},t)$ satisfies:

\begin{equation}
    \partial_t p(\mathbf{x},t) + \nabla \cdot (p(\mathbf{x},t) \mathbf{v}(\mathbf{x},t)) = 0.
\end{equation}

This equation ensures conservation of probability mass during the flow transformation.

\subsection{Conditional Probability Paths}

Given a data distribution $p_{\text{data}}(\mathbf{x}_0)$ and base distribution $p_{\text{base}}(\mathbf{x})$, we define conditional paths $q_t(\mathbf{x}|\mathbf{x}_0)$ such that:
\begin{equation}
    p_t(\mathbf{x}) = \int p_{\text{data}}(\mathbf{x}_0) q_t(\mathbf{x}|\mathbf{x}_0) d\mathbf{x}_0, \quad t \in [0,1],
\end{equation}
with boundary conditions $p_0 = p_{\text{data}}$ and $p_1 = p_{\text{base}}$.

The conditional velocity field is defined as:
\begin{equation}
    \mathbf{u}_t(\mathbf{x}|\mathbf{x}_0) := \mathbb{E}\left[\frac{d}{dt} \mathbf{x}_t \bigg| \mathbf{x}_t = \mathbf{x}, \mathbf{x}_0\right],
\end{equation}
which provides analytic supervision for learning the marginal velocity field $\mathbf{v}(\mathbf{x},t)$.

\section{Gaussian Conditional Paths}

\subsection{Mathematical Derivation}

For Gaussian conditional paths of the form:
\begin{equation}
    q_t(\mathbf{x}|\mathbf{x}_0) = \mathcal{N}(\mathbf{x}; \boldsymbol{\mu}_t(\mathbf{x}_0), \sigma_t^2\mathbf{I}),
\end{equation}
where $\mathbf{x}_t = \boldsymbol{\mu}_t(\mathbf{x}_0) + \sigma_t \boldsymbol{\epsilon}$ with $\boldsymbol{\epsilon} \sim \mathcal{N}(0,\mathbf{I})$.

Taking the time derivative:
\begin{equation}
    \frac{d}{dt} \mathbf{x}_t = \dot{\boldsymbol{\mu}}_t(\mathbf{x}_0) + \dot{\sigma}_t \boldsymbol{\epsilon},
\end{equation}
where $\boldsymbol{\epsilon} = \frac{\mathbf{x}_t - \boldsymbol{\mu}_t(\mathbf{x}_0)}{\sigma_t}$.

The conditional expectation gives us the target velocity:
\begin{equation}
    \mathbf{u}_t(\mathbf{x}|\mathbf{x}_0) = \dot{\boldsymbol{\mu}}_t(\mathbf{x}_0) + \frac{\dot{\sigma}_t}{\sigma_t}(\mathbf{x} - \boldsymbol{\mu}_t(\mathbf{x}_0)).
\end{equation}

\subsection{Path Choices}

Different choices of $\boldsymbol{\mu}_t(\mathbf{x}_0)$ and $\sigma_t$ lead to different flow matching formulations:

\textbf{1. Linear bridge (optimal transport):}
\begin{itemize}
    \item $\boldsymbol{\mu}_t(\mathbf{x}_0) = (1-t)\mathbf{x}_0$, $\sigma_t = t\sigma_{\max}$
    \item Minimizes Wasserstein-2 distance
    \item Velocity: $\mathbf{u}_t(\mathbf{x}|\mathbf{x}_0) = \frac{\sigma_{\max}}{t}(\mathbf{x} - \mathbf{x}_0)$
\end{itemize}

\textbf{2. VP-like (DDPM-inspired):}
\begin{itemize}
    \item $\boldsymbol{\mu}_t(\mathbf{x}_0) = \alpha(t)\mathbf{x}_0$, $\sigma_t = \sqrt{1-\alpha^2(t)}$
    \item Preserves energy structure
    \item Velocity: $\mathbf{u}_t(\mathbf{x}|\mathbf{x}_0) = \dot{\alpha}(t)\mathbf{x}_0 + \frac{\dot{\sigma}_t}{\sigma_t}(\mathbf{x} - \alpha(t)\mathbf{x}_0)$
\end{itemize}

\textbf{3. VE-like (SMLD-inspired):}
\begin{itemize}
    \item $\boldsymbol{\mu}_t(\mathbf{x}_0) = \mathbf{x}_0$, $\sigma_t = \sigma(t)$ increasing
    \item Preserves data structure
    \item Velocity: $\mathbf{u}_t(\mathbf{x}|\mathbf{x}_0) = \frac{\dot{\sigma}(t)}{\sigma(t)}(\mathbf{x} - \mathbf{x}_0)$
\end{itemize}

\section{Training and Inference}

\subsection{Flow Matching Objective}

The flow matching training objective is:
\begin{equation}
    J_{\text{FM}} = \mathbb{E}\left[\|\mathbf{v}_\theta(\mathbf{x}_t, t) - \mathbf{u}_t(\mathbf{x}_t|\mathbf{x}_0)\|^2\right],
\end{equation}
where the expectation is over $t \sim \mathcal{U}[0,1]$, $\mathbf{x}_0 \sim p_{\text{data}}$, and $\mathbf{x}_t \sim q_t(\cdot|\mathbf{x}_0)$.

\subsection{Mathematical Optimality}

The population minimizer of the flow matching objective is the posterior-averaged velocity field:
\begin{equation}
    \mathbf{v}^*(\mathbf{x},t) = \mathbb{E}\left[\mathbf{u}_t(\mathbf{x}_t|\mathbf{x}_0) \bigg| \mathbf{x}_t = \mathbf{x}, t\right].
\end{equation}

This result shows that flow matching learns the optimal velocity field that respects the continuity equation.

\subsection{Inference}

The learned velocity field $\mathbf{v}_\theta(\mathbf{x},t)$ defines a deterministic ODE:
\begin{equation}
    \frac{d}{dt} \mathbf{x} = \mathbf{v}_\theta(\mathbf{x},t), \quad t \in [0,1]
\end{equation}

The generation process involves:
\begin{enumerate}
    \item Initialize $\mathbf{x}(1) \sim p_1$ (e.g., $\mathcal{N}(0,\mathbf{I})$)
    \item Integrate ODE backwards: $\frac{d}{dt} \mathbf{x} = \mathbf{v}_\theta(\mathbf{x},t)$, $t: 1 \to 0$
    \item Use standard ODE solvers (Euler, RK4, Dormand-Prince, etc.)
\end{enumerate}

\section{Connection to Diffusion Models}

\subsection{Probability Flow ODE}

For a forward SDE $d\mathbf{x} = f(\mathbf{x},t)dt + g(t)dB_t$, the probability flow ODE is:
\begin{equation}
    \frac{d}{dt} \mathbf{x} = f(\mathbf{x},t) - \frac{1}{2}g^2(t)\nabla_{\mathbf{x}} \log p(\mathbf{x},t).
\end{equation}

When flow matching uses Gaussian paths that match the SDE's marginal distributions $p_t$, the learned velocity field $\mathbf{v}_\theta$ approximates the probability flow drift:
\begin{equation}
    \mathbf{v}_\theta(\mathbf{x},t) \approx f(\mathbf{x},t) - \frac{1}{2}g^2(t)\nabla_{\mathbf{x}} \log p(\mathbf{x},t)
\end{equation}

This connection allows flow matching to achieve the same approximation as diffusion models without ever computing the score function during training.

\section{MeanFlow: Average Velocity Field}

\subsection{Motivation}

MeanFlow addresses the challenge of single-step generation while maintaining high fidelity. The key idea is to learn a ground-truth average velocity field that summarizes an entire time interval.

\subsection{Definition}

For any pair $r < t$ and point $\mathbf{z}_t$ along the flow path:
\begin{equation}
    \overline{\mathbf{u}}(\mathbf{z}_t, r, t) := \frac{1}{t-r} \int_r^t \mathbf{v}(\mathbf{z}_\tau, \tau) d\tau.
\end{equation}

\subsection{MeanFlow Identity}

The core relationship is:
\begin{equation}
    \overline{\mathbf{u}}(\mathbf{z}_t, r, t) = \mathbf{v}(\mathbf{z}_t, t) - (t-r)\frac{d}{dt}\overline{\mathbf{u}}(\mathbf{z}_t, r, t),
\end{equation}
where the total derivative expands to:
\begin{equation}
    \frac{d}{dt}\overline{\mathbf{u}}(\mathbf{z}_t, r, t) = \mathbf{v}(\mathbf{z}_t, t) \cdot \nabla_{\mathbf{z}}\overline{\mathbf{u}}(\mathbf{z}_t, r, t) + \frac{\partial}{\partial t}\overline{\mathbf{u}}(\mathbf{z}_t, r, t).
\end{equation}

\subsection{Training and Sampling}

The training objective uses the MeanFlow identity to form targets:
\begin{equation}
    \mathcal{L}(\theta) = \mathbb{E}\left[\|\overline{\mathbf{u}}_\theta(\mathbf{z}_t, r, t) - \text{sg}(\overline{\mathbf{u}}_{\text{tgt}})\|_2^2\right],
\end{equation}
where $\text{sg}(\cdot)$ is the stop-gradient operator.

For 1-NFE sampling:
\begin{equation}
    \mathbf{x} = \mathbf{z}_0 = \boldsymbol{\epsilon} - \overline{\mathbf{u}}_\theta(\boldsymbol{\epsilon}, 0, 1)
\end{equation}

\section{Recent Advances and Applications}

\subsection{Practical Considerations}

Flow matching has several advantages over traditional score-based methods:
\begin{itemize}
    \item No score estimation required during training
    \item Analytic supervision through conditional paths
    \item Deterministic sampling with ODE integration
    \item Direct connection to optimal transport theory
\end{itemize}

\subsection{Applications}

Flow matching has been successfully applied to:
\begin{itemize}
    \item High-resolution image generation
    \item Audio synthesis
    \item Molecular generation
    \item Scientific computing applications
\end{itemize}

\section{Conclusion}

Flow matching provides a mathematically elegant and practically effective approach to generative modeling. The key insights include:

\begin{itemize}
    \item The continuity equation provides the theoretical foundation
    \item Conditional paths enable analytic supervision
    \item The connection to diffusion models reveals the underlying unity
    \item MeanFlow enables efficient single-step generation
\end{itemize}

The mathematical rigor and practical effectiveness of flow matching make it a valuable tool in the arsenal of generative modeling techniques, with ongoing research exploring its applications to increasingly complex domains.

\chapter{Variational Autoencoders: Theory and Applications}

\section{Introduction}

Variational Autoencoders (VAEs) represent a fundamental breakthrough in generative modeling that combines the power of neural networks with principled probabilistic inference. Unlike traditional approaches that require intractable likelihood computations, VAEs provide a scalable framework for learning complex data distributions through variational inference.

The key insight of VAEs is to approximate intractable posterior distributions using neural networks, enabling end-to-end training through the reparameterization trick. This chapter provides a comprehensive mathematical treatment of VAEs, from theoretical foundations to practical implementations and recent advances.

\section{The Challenge of Generative Modeling}

\subsection{Problem Formulation}

Given a dataset $\mathcal{D} = \{x^{(i)}\}_{i=1}^N$ sampled from an unknown distribution $p^*(x)$, we want to:
\begin{itemize}
    \item Learn a model $p_\theta(x)$ that approximates $p^*(x)$
    \item Generate new samples from the learned distribution
    \item Perform inference: $p_\theta(z|x)$ for latent variables $z$
\end{itemize}

\subsection{Traditional Approaches and Limitations}

\textbf{Maximum Likelihood Estimation (MLE):}
\begin{equation}
    \theta^* = \arg\max_\theta \sum_{i=1}^N \log p_\theta(x^{(i)})
\end{equation}

\textbf{Problem:} Direct likelihood computation is often intractable for complex models.

\textbf{Latent Variable Models:}
\begin{equation}
    p_\theta(x) = \int p_\theta(x|z) p_\theta(z) dz
\end{equation}

\textbf{Problem:} Integration over latent space $z$ is typically intractable.

\section{Variational Inference Framework}

\subsection{The ELBO}

Given a latent variable model $p_\theta(x,z) = p_\theta(x|z)p_\theta(z)$, we want to approximate the intractable posterior $p_\theta(z|x)$ and maximize the marginal likelihood $p_\theta(x)$.

The KL divergence can be decomposed as:
\begin{equation}
    \text{KL}(q_\phi(z|x) \| p_\theta(z|x)) = \log p_\theta(x) - \text{ELBO}(x; \theta, \phi)
\end{equation}

where the Evidence Lower Bound is:
\begin{equation}
    \text{ELBO}(x; \theta, \phi) = \mathbb{E}_{q_\phi(z|x)}[\log p_\theta(x|z)] - \text{KL}(q_\phi(z|x) \| p_\theta(z))
\end{equation}

Since $\text{KL} \geq 0$, we have:
\begin{equation}
    \log p_\theta(x) \geq \text{ELBO}(x; \theta, \phi)
\end{equation}

\subsection{ELBO Interpretation}

The ELBO consists of two components:

\textbf{1. Reconstruction term:} $\mathbb{E}_{q_\phi(z|x)}[\log p_\theta(x|z)]$
\begin{itemize}
    \item Encourages $q_\phi(z|x)$ to encode information about $x$
    \item Maximizes likelihood of reconstructing $x$ from $z$
\end{itemize}

\textbf{2. Regularization term:} $-\text{KL}(q_\phi(z|x) \| p_\theta(z))$
\begin{itemize}
    \item Encourages $q_\phi(z|x)$ to be close to prior $p_\theta(z)$
    \item Prevents overfitting to training data
\end{itemize}

\section{VAE Architecture and Training}

\subsection{Architecture}

\textbf{Encoder (Recognition model):} $q_\phi(z|x)$
\begin{itemize}
    \item Maps data $x$ to latent representation $z$
    \item Typically: $q_\phi(z|x) = \mathcal{N}(z; \mu_\phi(x), \sigma_\phi^2(x)I)$
\end{itemize}

\textbf{Decoder (Generative model):} $p_\theta(x|z)$
\begin{itemize}
    \item Reconstructs data $x$ from latent $z$
    \item Typically: $p_\theta(x|z) = \mathcal{N}(x; \mu_\theta(z), \sigma_\theta^2 I)$ or Bernoulli
\end{itemize}

\textbf{Prior:} $p_\theta(z) = \mathcal{N}(z; 0, I)$

\subsection{The Reparameterization Trick}

\textbf{Problem:} How to backpropagate through stochastic sampling $z \sim q_\phi(z|x)$?

\textbf{Solution:} Reparameterization trick
\begin{equation}
    z = \mu_\phi(x) + \sigma_\phi(x) \odot \epsilon, \quad \epsilon \sim \mathcal{N}(0, I)
\end{equation}

\textbf{Benefits:}
\begin{itemize}
    \item Gradient flows through deterministic path: $x \to \mu_\phi(x), \sigma_\phi(x) \to z$
    \item Stochasticity moved to $\epsilon$ (independent of parameters)
    \item Enables end-to-end training with backpropagation
\end{itemize}

\subsection{Training Objective}

For a dataset $\mathcal{D} = \{x^{(i)}\}_{i=1}^N$, the VAE loss is:
\begin{equation}
    \mathcal{L}(\theta, \phi) = \sum_{i=1}^N \text{ELBO}(x^{(i)}; \theta, \phi)
\end{equation}

Expanding the ELBO:
\begin{equation}
    \mathcal{L}(\theta, \phi) = \sum_{i=1}^N \left[ \mathbb{E}_{q_\phi(z|x^{(i)})}[\log p_\theta(x^{(i)}|z)] - \text{KL}(q_\phi(z|x^{(i)}) \| p_\theta(z)) \right]
\end{equation}

\subsection{Analytical KL Divergence}

For Gaussian distributions, the KL divergence has a closed form:
\begin{equation}
    \text{KL}(\mathcal{N}(\mu, \sigma^2) \| \mathcal{N}(0, I)) = \frac{1}{2} \sum_{j=1}^J (\mu_j^2 + \sigma_j^2 - \log \sigma_j^2 - 1)
\end{equation}

where $J$ is the dimension of the latent space.

\textbf{Interpretation:}
\begin{itemize}
    \item $\mu_j^2$: Penalizes large mean values
    \item $\sigma_j^2 - \log \sigma_j^2 - 1$: Encourages $\sigma_j \approx 1$
    \item Total effect: Regularizes latent space to be close to standard Gaussian
\end{itemize}

\section{VAE Variants and Extensions}

\subsection{Beta-VAE: Controlling Regularization}

\textbf{Motivation:} Balance between reconstruction and regularization

\textbf{$\beta$-VAE objective:}
\begin{equation}
    \mathcal{L}(\theta, \phi) = \mathbb{E}_{q_\phi(z|x)}[\log p_\theta(x|z)] - \beta \text{KL}(q_\phi(z|x) \| p_\theta(z))
\end{equation}

\textbf{Effect of $\beta$:}
\begin{itemize}
    \item $\beta > 1$: Stronger regularization $\Rightarrow$ Better disentanglement
    \item $\beta < 1$: Weaker regularization $\Rightarrow$ Better reconstruction
    \item $\beta = 1$: Standard VAE
\end{itemize}

\subsection{VQ-VAE: Vector Quantized VAE}

\textbf{Key idea:} Use discrete latent representations instead of continuous

\textbf{Architecture:}
\begin{itemize}
    \item Encoder: $z_e = E(x)$ (continuous)
    \item Quantization: $z_q = \text{Quantize}(z_e)$ (discrete)
    \item Decoder: $\hat{x} = D(z_q)$
\end{itemize}

\textbf{Quantization process:}
\begin{equation}
    z_q = \arg\min_{k} \|z_e - e_k\|^2
\end{equation}
where $\{e_k\}_{k=1}^K$ are learnable codebook vectors.

\textbf{Training objective:}
\begin{equation}
    \mathcal{L} = \log p(x|z_q) + \|\text{sg}(z_e) - e_k\|^2 + \beta\|z_e - \text{sg}(e_k)\|^2
\end{equation}

where $\text{sg}(\cdot)$ is stop-gradient operator.

\subsection{WAE: Wasserstein Autoencoder}

\textbf{Motivation:} Replace KL divergence with Wasserstein distance

\textbf{WAE objective:}
\begin{equation}
    \mathcal{L} = \mathbb{E}_{p(x)} \mathbb{E}_{q_\phi(z|x)} [c(x, D(z))] + \lambda \mathcal{D}(q_\phi(z), p(z))
\end{equation}

where:
\begin{itemize}
    \item $c(x, D(z))$: Reconstruction cost
    \item $\mathcal{D}(q_\phi(z), p(z))$: Distance between aggregated posterior and prior
    \item $\lambda$: Regularization weight
\end{itemize}

\subsection{VAE-GAN: Combining VAE and GAN}

\textbf{Idea:} Use GAN discriminator to improve VAE reconstruction quality

\textbf{Training objective:}
\begin{equation}
    \mathcal{L} = \mathcal{L}_{\text{VAE}} + \lambda \mathcal{L}_{\text{GAN}}
\end{equation}

where:
\begin{itemize}
    \item $\mathcal{L}_{\text{VAE}}$: Standard VAE loss
    \item $\mathcal{L}_{\text{GAN}}$: Adversarial loss for discriminator
\end{itemize}

\section{Advanced Topics}

\subsection{Hierarchical VAE}

\textbf{Motivation:} Learn hierarchical latent representations

\textbf{Architecture:}
\begin{equation}
    p_\theta(x, z_1, \ldots, z_L) = p_\theta(x|z_1) \prod_{l=1}^{L-1} p_\theta(z_l|z_{l+1}) p_\theta(z_L)
\end{equation}

\textbf{Inference:}
\begin{equation}
    q_\phi(z_{1:L}|x) = \prod_{l=1}^L q_\phi(z_l|z_{<l}, x)
\end{equation}

\subsection{Conditional VAE (C-VAE)}

\textbf{Extension:} Incorporate conditioning information

\textbf{Model:}
\begin{equation}
    p_\theta(x, z|c) = p_\theta(x|z, c) p_\theta(z)
\end{equation}

\textbf{Inference:}
\begin{equation}
    q_\phi(z|x, c) = \mathcal{N}(z; \mu_\phi(x, c), \sigma_\phi^2(x, c))
\end{equation}

\textbf{Training:}
\begin{equation}
    \mathcal{L} = \mathbb{E}_{q_\phi(z|x,c)}[\log p_\theta(x|z,c)] - \text{KL}(q_\phi(z|x,c) \| p_\theta(z))
\end{equation}

\subsection{Importance Weighted Autoencoder (IWAE)}

\textbf{Motivation:} Better approximation of log-likelihood

\textbf{IWAE objective:}
\begin{equation}
    \mathcal{L}_K = \mathbb{E}_{z_1, \ldots, z_K \sim q_\phi(z|x)} \left[ \log \frac{1}{K} \sum_{k=1}^K \frac{p_\theta(x, z_k)}{q_\phi(z_k|x)} \right]
\end{equation}

\textbf{Key properties:}
\begin{itemize}
    \item $\mathcal{L}_K \geq \mathcal{L}_{K-1}$ (monotonic improvement)
    \item $\lim_{K \to \infty} \mathcal{L}_K = \log p_\theta(x)$ (exact likelihood)
    \item $K=1$ recovers standard VAE
\end{itemize}

\section{Implementation and Practical Considerations}

\subsection{Common Issues and Solutions}

\textbf{Posterior collapse:}
\begin{itemize}
    \item \textbf{Symptom:} $q_\phi(z|x) \approx p_\theta(z)$ (no information in latent)
    \item \textbf{Causes:} Strong decoder, weak encoder
    \item \textbf{Solutions:} KL annealing, free bits, stronger encoder
\end{itemize}

\textbf{Blurry reconstructions:}
\begin{itemize}
    \item \textbf{Cause:} Gaussian decoder assumption
    \item \textbf{Solutions:} VAE-GAN, better decoder architectures
\end{itemize}

\textbf{Mode collapse:}
\begin{itemize}
    \item \textbf{Symptom:} Limited diversity in generated samples
    \item \textbf{Solutions:} Better regularization, architectural improvements
\end{itemize}

\subsection{Training Tricks}

\begin{itemize}
    \item \textbf{KL annealing:} Gradually increase KL weight
    \item \textbf{Free bits:} Prevent posterior collapse
    \item \textbf{Warm-up:} Start with reconstruction only
\end{itemize}

\section{Recent Advances and Applications}

\subsection{Transformer-based VAEs}

\begin{itemize}
    \item VQ-VAE-2: Hierarchical discrete representations
    \item DALL-E: Text-to-image generation
    \item VQGAN: High-quality image generation
\end{itemize}

\subsection{Improved Training}

\begin{itemize}
    \item NVAE: Normalizing flow decoder
    \item VDVAE: Very deep VAE
    \item VDM: Diffusion-based decoder
\end{itemize}

\subsection{Applications}

\begin{itemize}
    \item Text generation (GPT-style models)
    \item Image synthesis and editing
    \item Audio generation
    \item Scientific modeling
\end{itemize}

\section{Evaluation Metrics}

\textbf{Reconstruction quality:}
\begin{itemize}
    \item MSE, SSIM, FID (Fréchet Inception Distance)
    \item Perceptual losses
\end{itemize}

\textbf{Generation quality:}
\begin{itemize}
    \item FID, IS (Inception Score)
    \item Human evaluation
\end{itemize}

\textbf{Representation quality:}
\begin{itemize}
    \item Disentanglement metrics ($\beta$-VAE metric, MIG)
    \item Downstream task performance
    \item Interpolation smoothness
\end{itemize}

\section{Open Questions and Future Directions}

\subsection{Theoretical Understanding}

\begin{itemize}
    \item Why do VAEs work despite approximation errors?
    \item Optimal choice of prior distributions
    \item Connection to other generative models
\end{itemize}

\subsection{Technical Challenges}

\begin{itemize}
    \item Better posterior approximations
    \item Scalable training for large models
    \item Integration with other architectures
\end{itemize}

\subsection{Applications}

\begin{itemize}
    \item Multi-modal learning
    \item Scientific discovery
    \item Efficient compression
    \item Robust representations
\end{itemize}

\section{Conclusion}

VAEs have made fundamental contributions to generative modeling:

\textbf{Key contributions:}
\begin{itemize}
    \item Scalable variational inference for latent variable models
    \item End-to-end training with reparameterization trick
    \item Foundation for many generative models
\end{itemize}

\textbf{VAE variants address:}
\begin{itemize}
    \item Disentanglement ($\beta$-VAE)
    \item Discrete representations (VQ-VAE)
    \item Sample quality (VAE-GAN, WAE)
    \item Hierarchical modeling (HVAE)
\end{itemize}

\textbf{Impact:}
\begin{itemize}
    \item Influenced modern generative models
    \item Enabled practical latent variable modeling
    \item Foundation for many applications
\end{itemize}

The mathematical rigor and practical effectiveness of VAEs ensure their continued relevance in the rapidly evolving field of generative modeling.

\chapter{Normalizing Flows: Exact Likelihood Generative Modeling}

\section{Introduction}

Normalizing flows represent a powerful class of generative models that provide exact likelihood computation and exact sampling capabilities. Unlike other generative approaches that rely on approximations, normalizing flows learn invertible transformations that map simple distributions to complex data distributions while maintaining tractable likelihood evaluation.

The fundamental principle underlying normalizing flows is the change of variables theorem, which allows us to transform probability densities through invertible mappings. This chapter provides a comprehensive mathematical treatment of normalizing flows, from theoretical foundations to practical implementations and recent advances.

\section{Mathematical Foundations}

\subsection{Change of Variables Theorem}

Given a bijective transformation $T: \mathbb{R}^d \to \mathbb{R}^d$ with inverse $T^{-1}$, and probability densities $p_X(x)$ and $p_Y(y)$, we have:
\begin{equation}
    p_Y(y) = p_X(T^{-1}(y)) \left|\det \frac{\partial T^{-1}(y)}{\partial y}\right|
\end{equation}
where $\frac{\partial T^{-1}(y)}{\partial y}$ is the Jacobian matrix of the inverse transformation.

\textbf{Key insights:}
\begin{itemize}
    \item We can transform simple distributions into complex ones
    \item Computing determinants is computationally expensive
    \item We need to design transformations with tractable Jacobians
\end{itemize}

\subsection{Proof of Change of Variables Formula}

\textbf{Theorem:} For bijective $T: \mathbb{R}^d \to \mathbb{R}^d$ with inverse $T^{-1}$,
\begin{equation}
    p_Y(y) = p_X(T^{-1}(y)) \left|\det \frac{\partial T^{-1}(y)}{\partial y}\right|
\end{equation}

\textbf{Proof:} Let $A \subset \mathbb{R}^d$ be any measurable set. Then:
\begin{equation}
    \int_A p_Y(y) dy = \mathbb{P}(Y \in A) = \mathbb{P}(X \in T^{-1}(A))
\end{equation}

Using the change of variables formula for integrals:
\begin{equation}
    \mathbb{P}(X \in T^{-1}(A)) = \int_{T^{-1}(A)} p_X(x) dx = \int_A p_X(T^{-1}(y)) \left|\det \frac{\partial T^{-1}(y)}{\partial y}\right| dy
\end{equation}

Since this holds for any measurable set $A$, we have the desired result.

\subsection{Normalizing Flows}

A normalizing flow is a sequence of invertible transformations:
\begin{equation}
    \mathbf{x}_0 \xrightarrow{T_1} \mathbf{x}_1 \xrightarrow{T_2} \mathbf{x}_2 \xrightarrow{T_3} \cdots \xrightarrow{T_K} \mathbf{x}_K
\end{equation}

The probability density transforms as:
\begin{equation}
    \log p(\mathbf{x}_K) = \log p(\mathbf{x}_0) + \sum_{k=1}^K \log \left|\det \frac{\partial T_k}{\partial \mathbf{x}_{k-1}}\right|
\end{equation}

\textbf{Goal:} Learn $T_1, T_2, \ldots, T_K$ such that $p(\mathbf{x}_K)$ matches the data distribution.

\section{Autoregressive Flows}

\subsection{Autoregressive Transformations}

An autoregressive transformation processes variables in order:
\begin{equation}
    y_i = \tau(x_i; \mathbf{h}_i)
\end{equation}
where $\mathbf{h}_i = h_i(x_1, x_2, \ldots, x_{i-1})$ is a function of previous variables.

\textbf{Key properties:}
\begin{itemize}
    \item \textbf{Triangular Jacobian:} $\frac{\partial y_i}{\partial x_j} = 0$ for $j > i$
    \item \textbf{Efficient determinant:} $\det J = \prod_{i=1}^d \frac{\partial y_i}{\partial x_i}$
    \item \textbf{Easy inverse:} Can be computed sequentially
\end{itemize}

\subsection{Masked Autoregressive Flow (MAF)}

MAF uses neural networks with masking to ensure autoregressive structure:
\begin{equation}
    y_i = x_i \odot \exp(\alpha_i) + \beta_i
\end{equation}
where $\alpha_i, \beta_i$ are outputs of a masked neural network.

\textbf{Masking strategy:}
\begin{itemize}
    \item Input mask: Only allow connections from $x_1, \ldots, x_{i-1}$ to $\alpha_i, \beta_i$
    \item Ensures $\frac{\partial y_i}{\partial x_j} = 0$ for $j \geq i$
    \item Jacobian is lower triangular by construction
\end{itemize}

\subsection{Inverse Autoregressive Flow (IAF)}

IAF is the inverse of MAF:
\begin{equation}
    x_i = (y_i - \beta_i) \odot \exp(-\alpha_i)
\end{equation}

\textbf{Key differences from MAF:}
\begin{itemize}
    \item \textbf{Fast sampling:} Can generate samples in parallel
    \item \textbf{Slow density evaluation:} Requires sequential computation
    \item \textbf{Fast training:} When used for variational inference
\end{itemize}

\textbf{Training strategy:} Use IAF for fast sampling, MAF for fast density evaluation.

\subsection{TarFlow: Transformer-based Autoregressive Flow}

TarFlow is a Transformer-based variant of Masked Autoregressive Flows designed for high-resolution image synthesis.

\textbf{Key Features:}
\begin{itemize}
    \item \textbf{Autoregressive Transformer Blocks:} Stack of autoregressive Transformer blocks on image patches
    \item \textbf{Alternating Autoregression:} Alternates autoregression direction between layers
    \item \textbf{End-to-End Training:} Direct modeling and generation of pixels
    \item \textbf{State-of-the-art Performance:} First standalone normalizing flow to match diffusion model quality
\end{itemize}

\textbf{Architecture:}
\begin{itemize}
    \item \textbf{Input Processing:} Images divided into patches, each patch embedded into $D$-dim\break vectors
    \item \textbf{Transformer Layers:} $T$ autoregressive transformer layers with causal attention
    \item \textbf{Permutation Strategy:} Alternating permutations between layers for bidirectional modeling
    \item \textbf{Output:} Final latent representation $z^T$ in standard Gaussian space
\end{itemize}

\textbf{Forward Transformation:}
\begin{align}
    \tilde{z}^t &= \pi^t(z^t) \\
    z_i^{t+1} &= \begin{cases}
        \tilde{z}_i^t & i = 0 \\
        (\tilde{z}_i^t - \mu_i^t(\tilde{z}_{<i}^t)) \odot \exp(-\alpha_i^t(\tilde{z}_{<i}^t)) & i > 0
    \end{cases}
\end{align}

\textbf{Training Loss:}
\begin{equation}
    \min_f 0.5\|z^T\|_2^2 + \sum_{t=0}^{T-1} \sum_{i=1}^{N-1} \sum_{j=0}^{D-1} \alpha_i^t(z^t_{<i})_j
\end{equation}

\section{Coupling Layers}

\subsection{Affine Coupling Layers}

Split the input $\mathbf{x}$ into two parts: $\mathbf{x}_A$ and $\mathbf{x}_B$.

\textbf{Forward transformation:}
\begin{align}
    \mathbf{y}_A &= \mathbf{x}_A \\
    \mathbf{y}_B &= \mathbf{x}_B \odot \exp(\alpha(\mathbf{x}_A)) + \beta(\mathbf{x}_A)
\end{align}

\textbf{Jacobian:}
\begin{equation}
    J = \begin{pmatrix}
        I & 0 \\
        \frac{\partial \mathbf{y}_B}{\partial \mathbf{x}_A} & \text{diag}(\exp(\alpha(\mathbf{x}_A)))
    \end{pmatrix}
\end{equation}

\textbf{Determinant:} $\det J = \prod_{i \in B} \exp(\alpha_i(\mathbf{x}_A))$

\subsection{Real NVP: Real-valued Non-volume Preserving}

Real NVP uses alternating coupling layers:
\begin{enumerate}
    \item \textbf{Checkerboard masking:} Alternate which dimensions are transformed
    \item \textbf{Channel-wise masking:} For images, alternate RGB channels
\end{enumerate}

\textbf{Architecture:}
\begin{itemize}
    \item ResNet blocks for $\alpha$ and $\beta$ networks
    \item Batch normalization between layers
    \item Multi-scale architecture for images
\end{itemize}

\textbf{Advantages:}
\begin{itemize}
    \item \textbf{Exact likelihood:} Can compute $p(x)$ exactly
    \item \textbf{Exact sampling:} Can generate samples exactly
    \item \textbf{Stable training:} No gradient issues
\end{itemize}

\section{Continuous Normalizing Flows}

\subsection{Neural Ordinary Differential Equations}

Instead of discrete transformations, use continuous dynamics:
\begin{equation}
    \frac{d\mathbf{x}(t)}{dt} = f_\theta(\mathbf{x}(t), t)
\end{equation}

\textbf{Change of variables formula:}
\begin{equation}
    \log p(\mathbf{x}(t_1)) = \log p(\mathbf{x}(t_0)) + \int_{t_0}^{t_1} \text{tr}\left(\frac{\partial f_\theta}{\partial \mathbf{x}}\right) dt
\end{equation}

\textbf{Advantages:}
\begin{itemize}
    \item \textbf{Flexible architecture:} Any neural network can be used
    \item \textbf{Adaptive computation:} Can adjust number of function evaluations
    \item \textbf{Memory efficient:} No need to store intermediate states
\end{itemize}

\subsection{Proof of Neural ODE Change of Variables Formula}

\textbf{Theorem:} For the ODE $\frac{d\mathbf{x}(t)}{dt} = f_\theta(\mathbf{x}(t), t)$, the change of variables formula is:
\begin{equation}
    \log p(\mathbf{x}(t_1)) = \log p(\mathbf{x}(t_0)) + \int_{t_0}^{t_1} \text{tr}\left(\frac{\partial f_\theta}{\partial \mathbf{x}}\right) dt
\end{equation}

\textbf{Proof:} The probability density $p(\mathbf{x}, t)$ satisfies the continuity equation:
\begin{equation}
    \frac{\partial p(\mathbf{x}, t)}{\partial t} + \nabla \cdot (f_\theta(\mathbf{x}, t) p(\mathbf{x}, t)) = 0
\end{equation}

Expanding the divergence:
\begin{equation}
    \frac{\partial p(\mathbf{x}, t)}{\partial t} + \nabla \cdot (f_\theta(\mathbf{x}, t)) p(\mathbf{x}, t) + f_\theta(\mathbf{x}, t) \cdot \nabla p(\mathbf{x}, t) = 0
\end{equation}

Taking the log-derivative and substituting the ODE:
\begin{equation}
    \frac{d \log p(\mathbf{x}(t), t)}{dt} = -\nabla \cdot f_\theta(\mathbf{x}(t), t)
\end{equation}

Integrating from $t_0$ to $t_1$ and using the trace identity $\nabla \cdot f = \text{tr}\left(\frac{\partial f}{\partial \mathbf{x}}\right)$:
\begin{equation}
    \log p(\mathbf{x}(t_1)) = \log p(\mathbf{x}(t_0)) + \int_{t_0}^{t_1} \text{tr}\left(\frac{\partial f_\theta}{\partial \mathbf{x}}\right) dt
\end{equation}

\subsection{FFJORD: Free-form Jacobian of Reversible Dynamics}

FFJORD uses the continuous change of variables formula with Hutchinson's trace estimator:
\begin{equation}
    \text{tr}(A) = \mathbb{E}_{\mathbf{v}}[\mathbf{v}^T A \mathbf{v}]
\end{equation}
where $\mathbf{v}$ is a random vector with $\mathbb{E}[\mathbf{v}\mathbf{v}^T] = I$.

\textbf{Benefits:}
\begin{itemize}
    \item \textbf{Scalable:} Works for high-dimensional data
    \item \textbf{Flexible:} No architectural constraints
    \item \textbf{Exact:} No approximation errors
\end{itemize}

\section{Conditional Normalizing Flows}

\subsection{Goal and Conditional Likelihood}

\textbf{Problem:} Learn a high-dimensional conditional density $p_{\mathbf{Y}|\mathbf{X}}(\mathbf{y}|\mathbf{x})$ that is multi-modal and correlated across dimensions.

\textbf{CNF model:} Invertible, $\mathbf{x}$-conditioned map
\begin{equation}
    \mathbf{z} = f_\phi(\mathbf{y}; \mathbf{x}), \qquad f_\phi:\mathcal{Y}\times \mathcal{X}\to \mathcal{Z}
\end{equation}
with a simple conditional base $p_{\mathbf{Z}|\mathbf{X}}(\mathbf{z}|\mathbf{x})$ (e.g., diagonal Gaussian).

\textbf{Change of variables (conditional):}
\begin{equation}
    p_{\mathbf{Y}|\mathbf{X}}(\mathbf{y}|\mathbf{x}) = p_{\mathbf{Z}|\mathbf{X}}(f_\phi(\mathbf{y};\mathbf{x})|\mathbf{x}) \left|\det \frac{\partial f_\phi(\mathbf{y};\mathbf{x})}{\partial \mathbf{y}}\right|
\end{equation}

\textbf{Training:}
\begin{equation}
    \max_\phi \mathbb{E}_{(\mathbf{x},\mathbf{y})}\left[\log p_{\mathbf{Z}|\mathbf{X}}(f_\phi(\mathbf{y};\mathbf{x})|\mathbf{x}) + \log \left| \det \frac{\partial f_\phi(\mathbf{y};\mathbf{x})}{\partial \mathbf{y}} \right|\right]
\end{equation}

\textbf{Sampling:} $\mathbf{z}\sim p_{\mathbf{Z}|\mathbf{X}}(\cdot|\mathbf{x})$, then $\mathbf{y}= f_\phi^{-1}(\mathbf{z};\mathbf{x})$.

\subsection{Conditioning and Architecture}

\textbf{Condition injection:}
\begin{align}
    \text{Conditional prior: } & p(\mathbf{z}|\mathbf{x})=\mathcal{N}(\mu(\mathbf{x}), \text{diag}\sigma^2(\mathbf{x}))\\
    \text{Coupling (affine): } &
    \begin{cases}
        \mathbf{y}_a = \mathbf{z}_a, \\
        \mathbf{y}_b = \mathbf{z}_b \odot \exp s(\mathbf{z}_a,\mathbf{x}) + t(\mathbf{z}_a,\mathbf{x}),
    \end{cases}
    \Rightarrow \log|\det J|=\sum s(\mathbf{z}_a,\mathbf{x})
\end{align}

\textbf{Architecture primitives:}
\begin{itemize}
    \item \emph{Coupling layers} (additive/affine/rational-quadratic spline) for cheap Jacobians
    \item \emph{Invertible $1\times 1$ convs} or permutations to mix channels
    \item \emph{Multi-scale} (squeeze \& factor-out) for efficiency and expressivity
\end{itemize}

\section{Training and Applications}

\subsection{Maximum Likelihood Training}

Given data $\{\mathbf{x}^{(i)}\}_{i=1}^N$, maximize:
\begin{equation}
    \mathcal{L}(\theta) = \frac{1}{N} \sum_{i=1}^N \log p_\theta(\mathbf{x}^{(i)})
\end{equation}

\textbf{Gradient computation:}
\begin{equation}
    \nabla_\theta \mathcal{L} = \frac{1}{N} \sum_{i=1}^N \nabla_\theta \log p_\theta(\mathbf{x}^{(i)})
\end{equation}

\textbf{Challenges:}
\begin{itemize}
    \item \textbf{Mode collapse:} Flow might ignore some data modes
    \item \textbf{Training instability:} Gradients can be large
    \item \textbf{Computational cost:} Jacobian computation is expensive
\end{itemize}

\subsection{Applications}

\textbf{Generative modeling:}
\begin{itemize}
    \item Image generation (Real NVP, Glow)
    \item Audio synthesis (WaveGlow)
    \item Text generation (FlowSeq)
\end{itemize}

\textbf{Density estimation:}
\begin{itemize}
    \item Anomaly detection
    \item Out-of-distribution detection
    \item Uncertainty quantification
\end{itemize}

\textbf{Variational inference:}
\begin{itemize}
    \item Posterior approximation
    \item Bayesian neural networks
    \item Variational autoencoders
\end{itemize}

\section{Comparison with Other Methods}

\subsection{Normalizing Flows vs. Other Generative Models}

\begin{table}[h]
\centering
\begin{tabular}{|l|c|c|c|c|}
\hline
Method & Exact Likelihood & Fast Sampling & Fast Training & Stable \\
\hline
VAE & No & Yes & Yes & Yes \\
GAN & No & Yes & No & No \\
Diffusion & No & No & Yes & Yes \\
Flow & Yes & Yes & No & Yes \\
\hline
\end{tabular}
\caption{Comparison of generative modeling approaches}
\end{table}

\textbf{Unique advantages of flows:}
\begin{itemize}
    \item \textbf{Exact likelihood:} Can compute $p(x)$ exactly
    \item \textbf{Exact sampling:} Can generate samples exactly
    \item \textbf{Latent space:} Natural representation learning
    \item \textbf{Interpretability:} Understandable transformations
\end{itemize}

\section{Challenges and Future Directions}

\subsection{Current Challenges}

\textbf{Computational challenges:}
\begin{itemize}
    \item \textbf{Jacobian computation:} Expensive for high dimensions
    \item \textbf{Memory usage:} Need to store intermediate states
    \item \textbf{Training time:} Slower than GANs or VAEs
\end{itemize}

\textbf{Modeling challenges:}
\begin{itemize}
    \item \textbf{Architectural constraints:} Must be invertible
    \item \textbf{Mode collapse:} May ignore some data modes
    \item \textbf{Expressiveness:} Limited by transformation family
\end{itemize}

\textbf{Theoretical challenges:}
\begin{itemize}
    \item \textbf{Universal approximation:} Not guaranteed for all distributions
    \item \textbf{Convergence:} No global convergence guarantees
    \item \textbf{Generalization:} Limited theoretical understanding
\end{itemize}

\subsection{Future Directions}

\textbf{Architectural improvements:}
\begin{itemize}
    \item \textbf{More flexible transformations:} Beyond coupling layers
    \item \textbf{Better conditioning:} Improved numerical stability
    \item \textbf{Efficient computation:} Faster Jacobian computation
\end{itemize}

\textbf{Training improvements:}
\begin{itemize}
    \item \textbf{Better optimization:} More stable training
    \item \textbf{Regularization:} Prevent overfitting
    \item \textbf{Multi-scale training:} Progressive complexity
\end{itemize}

\textbf{Applications:}
\begin{itemize}
    \item \textbf{High-dimensional data:} Images, videos, audio
    \item \textbf{Scientific computing:} Physics simulations
    \item \textbf{Robotics:} Motion planning and control
\end{itemize}

\section{Conclusion}

Normalizing flows provide a unique approach to generative modeling with several key advantages:

\textbf{Key contributions:}
\begin{itemize}
    \item \textbf{Exact likelihood computation}
    \item \textbf{Exact sampling}
    \item \textbf{Flexible architectures}
    \item \textbf{Interpretable representations}
\end{itemize}

\textbf{Key insights:}
\begin{itemize}
    \item \textbf{Change of variables:} Foundation of all flow-based methods
    \item \textbf{Jacobian computation:} Central computational challenge
    \item \textbf{Architectural design:} Balance between expressiveness and efficiency
    \item \textbf{Training strategies:} Maximum likelihood with proper regularization
\end{itemize}

\textbf{Future work:}
\begin{itemize}
    \item \textbf{Scalability:} Handle higher-dimensional data
    \item \textbf{Efficiency:} Reduce computational costs
    \item \textbf{Theoretical understanding:} Better convergence guarantees
\end{itemize}

The mathematical rigor and practical effectiveness of normalizing flows make them a valuable tool in the arsenal of generative modeling techniques, with ongoing research exploring their applications to increasingly complex domains.

\chapter[Generative Adversarial Networks]{Generative Adversarial Networks:\\Adversarial Training for Generative Modeling}

\section{Introduction}

Generative Adversarial Networks (GANs) represent a revolutionary approach to generative modeling that introduced the concept of adversarial training. Unlike traditional methods that rely on likelihood maximization, GANs use a game-theoretic framework where two neural networks compete against each other: a generator that creates fake data and a discriminator that distinguishes between real and fake data.

The key insight of GANs is to frame generative modeling as a minimax game, where the generator learns to produce samples that are indistinguishable from real data, while the discriminator learns to accurately classify samples as real or fake. This chapter provides a comprehensive mathematical treatment of GANs, from theoretical foundations to practical implementations and recent advances.

\section{Mathematical Foundations}

\subsection{The Minimax Game}

The GAN objective is formulated as a minimax game:
\begin{equation}
    \min_G \max_D V(D, G) = \mathbb{E}_{x \sim p_{\text{data}}}[\log D(x)] + \mathbb{E}_{z \sim p_z}[\log(1 - D(G(z)))]
\end{equation}

where:
\begin{itemize}
    \item $G$ is the generator network that maps noise $z \sim p_z$ to data space
    \item $D$ is the discriminator network that outputs the probability that input $x$ is real
    \item $p_{\text{data}}$ is the real data distribution
    \item $p_z$ is the noise distribution (typically Gaussian)
\end{itemize}

\subsection{Optimal Discriminator}

For a fixed generator $G$, the optimal discriminator is:
\begin{equation}
    D^*(x) = \frac{p_{\text{data}}(x)}{p_{\text{data}}(x) + p_g(x)}
\end{equation}

where $p_g$ is the distribution induced by the generator.

\textbf{Proof:} The discriminator's objective is:
\begin{equation}
    \max_D V(D, G) = \int_x p_{\text{data}}(x) \log D(x) + p_g(x) \log(1 - D(x)) dx
\end{equation}

Taking the derivative with respect to $D(x)$ and setting it to zero:
\begin{equation}
    \frac{p_{\text{data}}(x)}{D(x)} - \frac{p_g(x)}{1 - D(x)} = 0
\end{equation}

Solving for $D(x)$ gives the optimal discriminator.

\subsection{Global Optimum}

When the generator is optimal, the global minimum of the minimax game is achieved when $p_g = p_{\text{data}}$, and the value of the game is $-\log 4$.

\textbf{Proof:} Substituting the optimal discriminator into the objective:
\begin{equation}
    C(G) = \max_D V(D, G) = \mathbb{E}_{x \sim p_{\text{data}}}[\log D^*(x)] + \mathbb{E}_{x \sim p_g}[\log(1 - D^*(x))]
\end{equation}

When $p_g = p_{\text{data}}$, we have $D^*(x) = \frac{1}{2}$ for all $x$, giving:
\begin{equation}
    C(G) = \mathbb{E}_{x \sim p_{\text{data}}}[\log \frac{1}{2}] + \mathbb{E}_{x \sim p_g}[\log \frac{1}{2}] = -\log 4
\end{equation}

\section{Training Dynamics and Challenges}

\subsection{Training Algorithm}

The GAN training algorithm alternates between:
\begin{enumerate}
    \item \textbf{Discriminator update:} $\theta_d \leftarrow \theta_d + \alpha_d \nabla_{\theta_d} V(D, G)$
    \item \textbf{Generator update:} $\theta_g \leftarrow \theta_g - \alpha_g \nabla_{\theta_g} V(D, G)$
\end{enumerate}

\subsection{Mode Collapse}

Mode collapse occurs when the generator produces a limited variety of samples, ignoring some modes of the data distribution.

\textbf{Causes:}
\begin{itemize}
    \item Generator finds a mode that consistently fools the discriminator
    \item Discriminator becomes too strong, causing generator to focus on specific samples
    \item Training instability leads to convergence to suboptimal solutions
\end{itemize}

\textbf{Solutions:}
\begin{itemize}
    \item Unrolled GANs: Use multiple discriminator updates
    \item Feature matching: Match intermediate representations
    \item Minibatch discrimination: Encourage diversity in generated samples
\end{itemize}

\subsection{Training Instability}

GAN training is notoriously unstable due to:
\begin{itemize}
    \item \textbf{Non-convex optimization:} The minimax game has no global convergence guarantees
    \item \textbf{Vanishing gradients:} When discriminator is too good, generator gradients vanish
    \item \textbf{Oscillatory behavior:} Generator and discriminator can oscillate without converging
\end{itemize}

\section{Improved GAN Architectures}

\subsection{DCGAN: Deep Convolutional GANs}

DCGAN introduced architectural guidelines for stable GAN training:

\textbf{Generator architecture:}
\begin{itemize}
    \item Use transposed convolutions for upsampling
    \item Use batch normalization in generator (except output layer)
    \item Use ReLU activations in generator
    \item Use tanh activation in output layer
\end{itemize}

\textbf{Discriminator architecture:}
\begin{itemize}
    \item Use strided convolutions for downsampling
    \item Use batch normalization in discriminator (except input layer)
    \item Use LeakyReLU activations in discriminator
\end{itemize}

\subsection{WGAN: Wasserstein GAN}

WGAN addresses training instability by using the Wasserstein distance instead of JS divergence.

\textbf{Wasserstein distance:}
\begin{equation}
    W(p_{\text{data}}, p_g) = \inf_{\gamma \in \Pi(p_{\text{data}}, p_g)} \mathbb{E}_{(x,y) \sim \gamma}[\|x - y\|]
\end{equation}

\textbf{WGAN objective:}
\begin{equation}
    \min_G \max_{D \in \mathcal{D}} \mathbb{E}_{x \sim p_{\text{data}}}[D(x)] - \mathbb{E}_{z \sim p_z}[D(G(z))]
\end{equation}

where $\mathcal{D}$ is the set of 1-Lipschitz functions.

\textbf{Benefits:}
\begin{itemize}
    \item More stable training
    \item Meaningful loss function (correlates with sample quality)
    \item Reduced mode collapse
\end{itemize}

\subsection{WGAN-GP: Gradient Penalty}

WGAN-GP enforces the Lipschitz constraint through gradient penalty:

\textbf{Objective:}
\begin{equation}
    \mathcal{L} = \mathbb{E}_{x \sim p_{\text{data}}}[D(x)] - \mathbb{E}_{z \sim p_z}[D(G(z))] + \lambda \mathbb{E}_{\hat{x} \sim p_{\hat{x}}}[(\|\nabla_{\hat{x}} D(\hat{x})\| - 1)^2]
\end{equation}

where $\hat{x} = \epsilon x + (1-\epsilon) G(z)$ with $\epsilon \sim U[0,1]$.

\section{Loss Functions and Training Techniques}

\subsection{Alternative Loss Functions}

\textbf{Least Squares GAN (LSGAN):}
\begin{equation}
    \min_D \frac{1}{2}\mathbb{E}_{x \sim p_{\text{data}}}[(D(x) - 1)^2] + \frac{1}{2}\mathbb{E}_{z \sim p_z}[(D(G(z)))^2]
\end{equation}

\textbf{Hinge Loss:}
\begin{equation}
    \mathcal{L}_D = -\mathbb{E}_{x \sim p_{\text{data}}}[\min(0, -1 + D(x))] - \mathbb{E}_{z \sim p_z}[\min(0, -1 - D(G(z)))]
\end{equation}

\subsection{Training Techniques}

\textbf{Spectral Normalization:}
\begin{itemize}
    \item Normalizes the weight matrices to have spectral norm 1
    \item Enforces Lipschitz constraint on discriminator
    \item More stable than gradient penalty
\end{itemize}

\textbf{Progressive Growing:}
\begin{itemize}
    \item Start with low-resolution images
    \item Gradually increase resolution during training
    \item Improves training stability and sample quality
\end{itemize}

\textbf{Self-Attention:}
\begin{itemize}
    \item Add self-attention layers to generator and discriminator
    \item Captures long-range dependencies in images
    \item Improves sample quality
\end{itemize}

\section{Evaluation Metrics}

\subsection{Inception Score (IS)}

The Inception Score measures both quality and diversity:
\begin{equation}
    \text{IS} = \exp(\mathbb{E}_{x \sim p_g}[D_{KL}(p(y|x) \| p(y))])
\end{equation}

where $p(y|x)$ is the conditional class distribution and $p(y)$ is the marginal class distribution.

\subsection{Fréchet Inception Distance (FID)}

FID measures the distance between real and generated distributions in feature space:
\begin{equation}
    \text{FID} = \|\mu_r - \mu_g\|^2 + \text{Tr}(\Sigma_r + \Sigma_g - 2(\Sigma_r \Sigma_g)^{1/2})
\end{equation}

where $(\mu_r, \Sigma_r)$ and $(\mu_g, \Sigma_g)$ are the mean and covariance of real and generated features.

\subsection{Precision and Recall}

\textbf{Precision:} Fraction of generated samples that are realistic
\textbf{Recall:} Fraction of real data modes that are covered by generated samples

\section{Recent Advances}

\subsection{StyleGAN: Style-Based Generator}

StyleGAN introduces a style-based generator architecture:

\textbf{Key innovations:}
\begin{itemize}
    \item \textbf{Style modulation:} Each layer is modulated by a style vector
    \item \textbf{Noise injection:} Add stochastic variation at each layer
    \item \textbf{Progressive growing:} Start with low resolution and grow
    \item \textbf{Mapping network:} Transform latent code to style vectors
\end{itemize}

\textbf{Architecture:}
\begin{equation}
    \mathbf{w} = f(\mathbf{z}), \quad \mathbf{y} = g(\mathbf{w}, \mathbf{n})
\end{equation}

where $f$ is the mapping network, $g$ is the synthesis network, and $\mathbf{n}$ is noise.

\subsection{BigGAN: Large Scale GAN Training}

BigGAN demonstrates that GANs can scale to large datasets:

\textbf{Key techniques:}
\begin{itemize}
    \item \textbf{Large batch sizes:} Use batch sizes of 2048 or higher
    \item \textbf{Class conditioning:} Use class labels for conditional generation
    \item \textbf{Truncation trick:} Truncate latent codes for better quality
    \item \textbf{Orthogonal initialization:} Initialize weights orthogonally
\end{itemize}

\subsection{Self-Attention GAN (SAGAN)}

SAGAN introduces self-attention mechanisms to GANs:

\textbf{Self-attention layer:}
\begin{equation}
    \text{Attention}(Q, K, V) = \text{softmax}\left(\frac{QK^T}{\sqrt{d_k}}\right)V
\end{equation}

\textbf{Benefits:}
\begin{itemize}
    \item Captures long-range dependencies
    \item Improves sample quality
    \item Enables generation of complex scenes
\end{itemize}

\section{Applications}

\subsection{Image Generation}

\begin{itemize}
    \item \textbf{Unconditional generation:} Generate images from noise
    \item \textbf{Conditional generation:} Generate images from class labels or text
    \item \textbf{Style transfer:} Transfer artistic style between images
    \item \textbf{Super-resolution:} Generate high-resolution images from low-resolution inputs
\end{itemize}

\subsection{Data Augmentation}

\begin{itemize}
    \item \textbf{Synthetic data:} Generate additional training data
    \item \textbf{Domain adaptation:} Adapt models to new domains
    \item \textbf{Privacy preservation:} Generate synthetic data for privacy-sensitive applications
\end{itemize}

\subsection{Scientific Applications}

\begin{itemize}
    \item \textbf{Molecular generation:} Generate new drug molecules
    \item \textbf{Protein design:} Design new protein structures
    \item \textbf{Material discovery:} Discover new materials with desired properties
\end{itemize}

\section{Challenges and Limitations}

\subsection{Training Challenges}

\begin{itemize}
    \item \textbf{Instability:} Training can be unstable and sensitive to hyperparameters
    \item \textbf{Mode collapse:} Generator may ignore some data modes
    \item \textbf{Evaluation:} Difficult to evaluate GAN performance objectively
    \item \textbf{Computational cost:} Training can be computationally expensive
\end{itemize}

\subsection{Theoretical Limitations}

\begin{itemize}
    \item \textbf{No convergence guarantees:} No global convergence guarantees
    \item \textbf{Mode collapse:} Theoretical understanding of mode collapse is limited
    \item \textbf{Generalization:} Limited understanding of generalization properties
\end{itemize}

\section{Future Directions}

\subsection{Theoretical Advances}

\begin{itemize}
    \item \textbf{Convergence analysis:} Better understanding of training dynamics
    \item \textbf{Generalization bounds:} Theoretical guarantees for GAN performance
    \item \textbf{Mode collapse:} Better understanding and prevention of mode collapse
\end{itemize}

\subsection{Architectural Improvements}

\begin{itemize}
    \item \textbf{More stable training:} New architectures that are easier to train
    \item \textbf{Better conditioning:} Improved conditional generation
    \item \textbf{Efficiency:} More efficient training and inference
\end{itemize}

\subsection{Applications}

\begin{itemize}
    \item \textbf{High-resolution generation:} Generate higher resolution images
    \item \textbf{Video generation:} Extend GANs to video generation
    \item \textbf{3D generation:} Generate 3D objects and scenes
\end{itemize}

\section{Conclusion}

Generative Adversarial Networks have revolutionized the field of generative modeling by introducing adversarial training:

\textbf{Key contributions:}
\begin{itemize}
    \item \textbf{Adversarial training:} Game-theoretic framework for generative modeling
    \item \textbf{High-quality samples:} Generate realistic images and other data
    \item \textbf{Flexible conditioning:} Support for various types of conditioning
    \item \textbf{Active research area:} Continues to drive innovation in generative modeling
\end{itemize}

\textbf{Key insights:}
\begin{itemize}
    \item \textbf{Minimax game:} Generator and discriminator compete in a zero-sum game
    \item \textbf{Training challenges:} Instability and mode collapse are major challenges
    \item \textbf{Evaluation difficulty:} Objective evaluation of GAN performance remains challenging
    \item \textbf{Practical impact:} GANs have enabled many practical applications
\end{itemize}

\textbf{Future work:}
\begin{itemize}
    \item \textbf{Theoretical understanding:} Better convergence and generalization guarantees
    \item \textbf{Training stability:} More robust training procedures
    \item \textbf{Evaluation metrics:} Better ways to evaluate GAN performance
    \item \textbf{Applications:} Extending GANs to new domains and tasks
\end{itemize}

The adversarial training paradigm introduced by GANs has fundamentally changed how we approach generative modeling, with ongoing research exploring new architectures, training techniques, and applications.


% Part II: Autoregressive Models
\part{Autoregressive Models}

\chapter{Autoregressive Models: Sequential Generation and Language Modeling}

\section{Introduction}

Autoregressive models represent a fundamental approach to generative modeling that generates data sequentially, one element at a time, by modeling the conditional probability of each element given the previous elements. This approach has been particularly successful in language modeling, where the goal is to predict the next word given the previous words in a sequence.

The key insight of autoregressive models is to factorize the joint probability distribution as a product of conditional distributions:
\begin{equation}
    p(x_1, x_2, \ldots, x_T) = \prod_{t=1}^T p(x_t | x_{<t})
\end{equation}

This chapter provides a comprehensive mathematical treatment of autoregressive models, from theoretical foundations to practical implementations and recent advances, with particular focus on transformer-based architectures.

\section{Mathematical Foundations}

\subsection{Autoregressive Factorization}

Given a sequence $\mathbf{x} = (x_1, x_2, \ldots, x_T)$, the autoregressive factorization is:
\begin{equation}
    p(\mathbf{x}) = \prod_{t=1}^T p(x_t | x_1, x_2, \ldots, x_{t-1}) = \prod_{t=1}^T p(x_t | \mathbf{x}_{<t})
\end{equation}

where $\mathbf{x}_{<t} = (x_1, x_2, \ldots, x_{t-1})$ represents the history up to time $t-1$.

\subsection{Maximum Likelihood Training}

The training objective is to maximize the log-likelihood:
\begin{equation}
    \mathcal{L}(\theta) = \sum_{t=1}^T \log p_\theta(x_t | \mathbf{x}_{<t})
\end{equation}

This can be optimized using standard gradient descent methods.

\subsection{Sampling}

To generate new sequences, we sample from the conditional distributions:
\begin{enumerate}
    \item Sample $x_1 \sim p_\theta(x_1)$
    \item For $t = 2, \ldots, T$: Sample $x_t \sim p_\theta(x_t | \mathbf{x}_{<t})$
\end{enumerate}

\section{Neural Autoregressive Models}

\subsection{Neural Language Models}

Neural autoregressive models use neural networks to parameterize the conditional distributions:
\begin{equation}
    p_\theta(x_t | \mathbf{x}_{<t}) = \text{softmax}(f_\theta(\mathbf{x}_{<t}))_{x_t}
\end{equation}

where $f_\theta$ is a neural network that maps the history to logits over the vocabulary.

\subsection{Recurrent Neural Networks (RNNs)}

RNNs process sequences sequentially using hidden states:
\begin{align}
    h_t &= f_\theta(h_{t-1}, x_t) \\
    p(x_{t+1} | \mathbf{x}_{\leq t}) &= \text{softmax}(W h_t + b)
\end{align}

\textbf{Advantages:}
\begin{itemize}
    \item Can handle variable-length sequences
    \item Memory of previous states
    \item Sequential processing is natural
\end{itemize}

\textbf{Disadvantages:}
\begin{itemize}
    \item Sequential processing is slow
    \item Vanishing/exploding gradients
    \item Limited long-range dependencies
\end{itemize}

\subsection{Long Short-Term Memory (LSTM)}

LSTM addresses the vanishing gradient problem through gating mechanisms:
\begin{align}
    f_t &= \sigma(W_f \cdot [h_{t-1}, x_t] + b_f) \quad \text{(forget gate)} \\
    i_t &= \sigma(W_i \cdot [h_{t-1}, x_t] + b_i) \quad \text{(input gate)} \\
    \tilde{C}_t &= \tanh(W_C \cdot [h_{t-1}, x_t] + b_C) \quad \text{(candidate values)} \\
    C_t &= f_t * C_{t-1} + i_t * \tilde{C}_t \quad \text{(cell state)} \\
    o_t &= \sigma(W_o \cdot [h_{t-1}, x_t] + b_o) \quad \text{(output gate)} \\
    h_t &= o_t * \tanh(C_t) \quad \text{(hidden state)}
\end{align}

\section{Transformer Architecture}

\subsection{Self-Attention Mechanism}

The core innovation of transformers is the self-attention mechanism:
\begin{equation}
    \text{Attention}(Q, K, V) = \text{softmax}\left(\frac{QK^T}{\sqrt{d_k}}\right)V
\end{equation}

where $Q$, $K$, and $V$ are query, key, and value matrices respectively.

\subsection{Multi-Head Attention}

Multi-head attention allows the model to attend to different representation subspaces:
\begin{equation}
    \text{MultiHead}(Q, K, V) = \text{Concat}(\text{head}_1, \ldots, \text{head}_h)W^O
\end{equation}

where:
\begin{equation}
    \text{head}_i = \text{Attention}(QW_i^Q, KW_i^K, VW_i^V)
\end{equation}

\subsection{Positional Encoding}

Since transformers have no inherent notion of sequence order, positional encodings are added:
\begin{align}
    PE_{(pos, 2i)} &= \sin(pos / 10000^{2i/d_{model}}) \\
    PE_{(pos, 2i+1)} &= \cos(pos / 10000^{2i/d_{model}})
\end{align}

\subsection{Transformer Block}

A transformer block consists of:
\begin{enumerate}
    \item Multi-head self-attention
    \item Residual connection and layer normalization
    \item Feed-forward network
    \item Another residual connection and layer normalization
\end{enumerate}

\section{GPT: Generative Pre-trained Transformer}

\subsection{Architecture}

GPT uses a decoder-only transformer architecture:
\begin{itemize}
    \item \textbf{Decoder-only:} No encoder, only decoder blocks
    \item \textbf{Causal attention:} Each position can only attend to previous positions
    \item \textbf{Autoregressive:} Generates text one token at a time
\end{itemize}

\subsection{Training}

GPT is trained using the standard autoregressive objective:
\begin{equation}
    \mathcal{L} = \sum_{t=1}^T \log p_\theta(x_t | \mathbf{x}_{<t})
\end{equation}

\subsection{Pre-training and Fine-tuning}

\textbf{Pre-training:} Train on large unlabeled text corpora
\textbf{Fine-tuning:} Adapt to specific tasks with labeled data

\section{BERT: Bidirectional Encoder Representations}

\subsection{Masked Language Modeling}

BERT uses masked language modeling (MLM) for pre-training:
\begin{itemize}
    \item Randomly mask 15\% of tokens
    \item Predict masked tokens from context
    \item Use bidirectional context (unlike GPT)
\end{itemize}

\subsection{Next Sentence Prediction}

BERT also uses next sentence prediction (NSP):
\begin{itemize}
    \item Given two sentences, predict if second follows first
    \item Helps with understanding sentence relationships
\end{itemize}

\section{Advanced Architectures}

\subsection{T5: Text-to-Text Transfer Transformer}

T5 frames all tasks as text-to-text problems:
\begin{itemize}
    \item \textbf{Unified framework:} All tasks use same input/output format
    \item \textbf{Encoder-decoder:} Uses both encoder and decoder
    \item \textbf{Transfer learning:} Pre-train on large dataset, fine-tune on specific tasks
\end{itemize}

\subsection{PaLM: Pathways Language Model}

PaLM uses the Pathways system for efficient training:
\begin{itemize}
    \item \textbf{Model parallelism:} Distribute model across multiple devices
    \item \textbf{Data parallelism:} Distribute data across multiple devices
    \item \textbf{Pipeline parallelism:} Overlap computation and communication
\end{itemize}

\section{Training Techniques}

\subsection{Gradient Accumulation}

For large models that don't fit in memory:
\begin{itemize}
    \item Accumulate gradients over multiple mini-batches
    \item Update parameters after accumulating sufficient gradients
    \item Effective batch size = mini-batch size × accumulation steps
\end{itemize}

\subsection{Mixed Precision Training}

Use both FP16 and FP32 precision:
\begin{itemize}
    \item \textbf{Forward pass:} Use FP16 for speed and memory
    \item \textbf{Gradient computation:} Use FP16 for gradients
    \item \textbf{Parameter updates:} Use FP32 for stability
\end{itemize}

\subsection{Gradient Clipping}

Prevent exploding gradients:
\begin{equation}
    \text{if } \|\mathbf{g}\| > \text{max\_norm: } \mathbf{g} \leftarrow \mathbf{g} \times \frac{\text{max\_norm}}{\|\mathbf{g}\|}
\end{equation}

\section{Evaluation Metrics}

\subsection{Perplexity}

Perplexity measures how well the model predicts the next token:
\begin{equation}
    \text{Perplexity} = \exp\left(-\frac{1}{T} \sum_{t=1}^T \log p_\theta(x_t | \mathbf{x}_{<t})\right)
\end{equation}

Lower perplexity indicates better performance.

\subsection{BLEU Score}

BLEU measures the quality of generated text by comparing n-grams:
\begin{equation}
    \text{BLEU} = \exp\left(\sum_{n=1}^N w_n \log p_n\right)
\end{equation}

where $p_n$ is the precision of n-grams and $w_n$ are weights.

\subsection{ROUGE Score}

ROUGE measures recall-oriented evaluation:
\begin{itemize}
    \item \textbf{ROUGE-N:} N-gram overlap
    \item \textbf{ROUGE-L:} Longest common subsequence
    \item \textbf{ROUGE-W:} Weighted longest common subsequence
\end{itemize}

\section{Applications}

\subsection{Language Modeling}

\begin{itemize}
    \item \textbf{Text generation:} Generate coherent and contextually appropriate text
    \item \textbf{Completion:} Complete partial sentences or paragraphs
    \item \textbf{Question answering:} Answer questions based on context
\end{itemize}

\subsection{Machine Translation}

\begin{itemize}
    \item \textbf{Sequence-to-sequence:} Translate between languages
    \item \textbf{Multilingual models:} Handle multiple languages in one model
    \item \textbf{Zero-shot translation:} Translate between language pairs not seen during training
\end{itemize}

\subsection{Text Summarization}

\begin{itemize}
    \item \textbf{Extractive:} Select important sentences from source
    \item \textbf{Abstractive:} Generate new sentences that capture main ideas
    \item \textbf{Multi-document:} Summarize information from multiple sources
\end{itemize}

\subsection{Dialog Systems}

\begin{itemize}
    \item \textbf{Conversational AI:} Generate human-like responses
    \item \textbf{Task-oriented:} Help users accomplish specific tasks
    \item \textbf{Open-domain:} Engage in general conversation
\end{itemize}

\section{Challenges and Limitations}

\subsection{Computational Challenges}

\begin{itemize}
    \item \textbf{Quadratic complexity:} Self-attention scales as $O(n^2)$ with sequence length
    \item \textbf{Memory requirements:} Large models require significant memory
    \item \textbf{Training time:} Long training times for large models
\end{itemize}

\subsection{Theoretical Limitations}

\begin{itemize}
    \item \textbf{Context length:} Limited by attention mechanism
    \item \textbf{Positional encoding:} May not generalize to longer sequences
    \item \textbf{Inductive bias:} Less inductive bias than RNNs
\end{itemize}

\subsection{Practical Challenges}

\begin{itemize}
    \item \textbf{Reproducibility:} Difficult to reproduce results due to computational requirements
    \item \textbf{Interpretability:} Black-box nature makes interpretation difficult
    \item \textbf{Bias:} Models may learn and amplify societal biases
\end{itemize}

\section{Recent Advances}

\subsection{Efficient Attention Mechanisms}

\textbf{Sparse Attention:}
\begin{itemize}
    \item \textbf{Local attention:} Only attend to nearby positions
    \item \textbf{Strided attention:} Attend to every k-th position
    \item \textbf{Random attention:} Randomly sample positions to attend to
\end{itemize}

\textbf{Linear Attention:}
\begin{itemize}
    \item \textbf{Performer:} Use random features to approximate attention
    \item \textbf{Linformer:} Low-rank approximation of attention matrix
    \item \textbf{Linear Transformer:} Linear complexity attention
\end{itemize}

\subsection{Long-Range Dependencies}

\textbf{Transformer-XL:}
\begin{itemize}
    \item \textbf{Segment-level recurrence:} Reuse hidden states across segments
    \item \textbf{Relative positional encoding:} Use relative positions instead of absolute
    \item \textbf{Longer context:} Handle longer sequences than standard transformers
\end{itemize}

\textbf{Compressive Transformer:}
\begin{itemize}
    \item \textbf{Compression:} Compress old memories to save space
    \item \textbf{Longer memory:} Maintain longer context than Transformer-XL
    \item \textbf{Efficient:} More efficient than standard attention
\end{itemize}

\section{Future Directions}

\subsection{Architectural Improvements}

\begin{itemize}
    \item \textbf{More efficient attention:} Reduce computational complexity
    \item \textbf{Better positional encoding:} Handle variable-length sequences
    \item \textbf{Hierarchical models:} Multi-scale processing
\end{itemize}

\subsection{Training Improvements}

\begin{itemize}
    \item \textbf{More efficient training:} Reduce computational requirements
    \item \textbf{Better optimization:} New optimization algorithms
    \item \textbf{Regularization:} Better regularization techniques
\end{itemize}

\subsection{Applications}

\begin{itemize}
    \item \textbf{Multimodal models:} Handle text, images, and other modalities
    \item \textbf{Code generation:} Generate and understand code
    \item \textbf{Scientific applications:} Apply to scientific problems
\end{itemize}

\section{Conclusion}

Autoregressive models have revolutionized natural language processing and generative modeling:

\textbf{Key contributions:}
\begin{itemize}
    \item \textbf{Sequential generation:} Generate data one element at a time
    \item \textbf{Transformer architecture:} Efficient parallel processing
    \item \textbf{Pre-training:} Learn from large unlabeled datasets
    \item \textbf{Transfer learning:} Adapt to specific tasks through fine-tuning
\end{itemize}

\textbf{Key insights:}
\begin{itemize}
    \item \textbf{Autoregressive factorization:} Factorize joint probability as product of conditionals
    \item \textbf{Attention mechanism:} Capture long-range dependencies efficiently
    \item \textbf{Pre-training importance:} Large-scale pre-training is crucial for performance
    \item \textbf{Scalability:} Performance improves with model size and data
\end{itemize}

\textbf{Future work:}
\begin{itemize}
    \item \textbf{Efficiency:} More efficient architectures and training methods
    \item \textbf{Longer context:} Handle longer sequences and documents
    \item \textbf{Multimodal:} Extend to multiple modalities
    \item \textbf{Interpretability:} Better understanding of model behavior
\end{itemize}

The autoregressive approach, combined with transformer architectures, has enabled unprecedented advances in natural language processing and continues to drive innovation in the field.


% Part III: Advanced Topics
\part{Advanced Topics}

\chapter{Discrete Diffusion Models: Extending Diffusion to Discrete Data}

\section{Introduction}

While traditional diffusion models work well with continuous data like images, many important applications involve discrete data such as text, categorical variables, or discrete tokens. Discrete diffusion models extend the diffusion framework to handle discrete data by defining appropriate forward and reverse processes on discrete spaces.

This chapter provides a comprehensive treatment of discrete diffusion models, covering theoretical foundations, practical implementations, and applications to various discrete data types.

\section{Mathematical Foundations}

\subsection{Discrete State Spaces}

For discrete data, we work with finite state spaces $\mathcal{X} = \{1, 2, \ldots, K\}$ where $K$ is the number of possible states. The forward process gradually adds noise to discrete data, while the reverse process learns to denoise and recover the original data.

\subsection{Forward Process}

The forward process is defined by a transition matrix $Q_t$ that specifies how to add noise at time $t$:
\begin{equation}
    q(x_t | x_{t-1}) = \text{Cat}(x_t; p = x_{t-1} Q_t)
\end{equation}

where $\text{Cat}$ denotes the categorical distribution.

\subsection{Reverse Process}

The reverse process learns to denoise by parameterizing the transition probabilities:
\begin{equation}
    p_\theta(x_{t-1} | x_t) = \text{Cat}(x_{t-1}; p = \text{softmax}(\log \hat{x}_{t-1} + \log Q_t^T))
\end{equation}

where $\hat{x}_{t-1}$ is the predicted logits from the neural network.

\section{Noise Schedules for Discrete Data}

\subsection{Uniform Noise}

The simplest approach is to add uniform noise:
\begin{equation}
    Q_t = (1 - \beta_t) I + \beta_t \frac{1}{K} \mathbf{1} \mathbf{1}^T
\end{equation}

where $\beta_t$ is the noise schedule and $\mathbf{1}$ is the all-ones vector.

\subsection{Gaussian-like Noise}

For discrete data that can be embedded in continuous space, we can use Gaussian noise in the embedding space:
\begin{equation}
    q(x_t | x_0) = \int q(x_t | \mathbf{z}_t) q(\mathbf{z}_t | \mathbf{z}_0) d\mathbf{z}_t
\end{equation}

where $\mathbf{z}_t$ is the continuous embedding of $x_t$.

\section{Training Objectives}

\subsection{Denoising Score Matching}

The training objective is to minimize the cross-entropy loss:
\begin{equation}
    \mathcal{L} = \mathbb{E}_{t, x_0, x_t} [-\log p_\theta(x_0 | x_t)]
\end{equation}

where $x_0$ is the original data and $x_t$ is the noisy version at time $t$.

\subsection{Simplified Objective}

For uniform noise, the objective simplifies to:
\begin{equation}
    \mathcal{L} = \mathbb{E}_{t, x_0, x_t} [-\log p_\theta(x_0 | x_t)]
\end{equation}

where $x_t$ is sampled from the forward process.

\section{Applications to Text}

\subsection{Text Diffusion}

For text data, we can apply discrete diffusion at the token level:
\begin{itemize}
    \item \textbf{Token-level diffusion:} Add noise to individual tokens
    \item \textbf{Character-level diffusion:} Add noise to individual characters
    \item \textbf{Subword-level diffusion:} Add noise to subword units
\end{itemize}

\subsection{Training Procedure}

\begin{enumerate}
    \item Sample a text sequence $x_0$
    \item Sample a time step $t \sim \text{Uniform}(0, T)$
    \item Sample noisy sequence $x_t$ from $q(x_t | x_0)$
    \item Train model to predict $x_0$ from $x_t$
\end{enumerate}

\subsection{Generation Procedure}

\begin{enumerate}
    \item Start with random noise $x_T \sim \text{Uniform}(1, K)$
    \item For $t = T, T-1, \ldots, 1$:
    \begin{itemize}
        \item Predict $\hat{x}_0$ from $x_t$
        \item Sample $x_{t-1}$ from $p_\theta(x_{t-1} | x_t)$
    \end{itemize}
\end{enumerate}

\section{Advanced Techniques}

\subsection{Learned Noise Schedules}

Instead of fixed noise schedules, we can learn the transition matrices:
\begin{equation}
    Q_t = \text{softmax}(W_t)
\end{equation}

where $W_t$ is a learnable matrix.

\subsection{Conditional Generation}

For conditional generation, we can condition on additional information:
\begin{equation}
    p_\theta(x_{t-1} | x_t, c) = \text{Cat}(x_{t-1}; p = \text{softmax}(\log \hat{x}_{t-1} + \log Q_t^T))
\end{equation}

where $c$ is the conditioning information.

\subsection{Classifier-Free Guidance}

Similar to continuous diffusion, we can use classifier-free guidance:
\begin{equation}
    \hat{x}_{t-1} = (1 + w) \hat{x}_{t-1}^c - w \hat{x}_{t-1}^u
\end{equation}

where $w$ is the guidance weight, and $c$ and $u$ denote conditional and unconditional predictions.

\section{Discrete Diffusion for Images}

\subsection{VQ-VAE Integration}

We can combine discrete diffusion with VQ-VAE:
\begin{enumerate}
    \item Encode images to discrete tokens using VQ-VAE
    \item Apply discrete diffusion to the tokens
    \item Decode the denoised tokens back to images
\end{enumerate}

\subsection{Benefits}

\begin{itemize}
    \item \textbf{Exact likelihood:} Can compute exact likelihood
    \item \textbf{Fast sampling:} Can use fewer steps than continuous diffusion
    \item \textbf{Better control:} More control over the generation process
\end{itemize}

\section{Evaluation Metrics}

\subsection{Perplexity}

For text generation, we can measure perplexity:
\begin{equation}
    \text{Perplexity} = \exp\left(-\frac{1}{N} \sum_{i=1}^N \log p(x_i)\right)
\end{equation}

\subsection{BLEU Score}

For machine translation tasks, we can use BLEU score:
\begin{equation}
    \text{BLEU} = \exp\left(\sum_{n=1}^N w_n \log p_n\right)
\end{equation}

\subsection{FID Score}

For image generation, we can use FID score:
\begin{equation}
    \text{FID} = \|\mu_r - \mu_g\|^2 + \text{Tr}(\Sigma_r + \Sigma_g - 2(\Sigma_r \Sigma_g)^{1/2})
\end{equation}

\section{Challenges and Limitations}

\subsection{Computational Challenges}

\begin{itemize}
    \item \textbf{Memory requirements:} Need to store transition matrices
    \item \textbf{Training time:} Can be slower than continuous diffusion
    \item \textbf{Sampling speed:} May require multiple steps for good quality
\end{itemize}

\subsection{Modeling Challenges}

\begin{itemize}
    \item \textbf{Discrete nature:} Harder to model than continuous data
    \item \textbf{State space size:} Large state spaces can be challenging
    \item \textbf{Noise design:} Choosing appropriate noise schedules is difficult
\end{itemize}

\section{Recent Advances}

\subsection{Improved Noise Schedules}

\begin{itemize}
    \item \textbf{Learned schedules:} Learn noise schedules from data
    \item \textbf{Adaptive schedules:} Adapt schedules based on data complexity
    \item \textbf{Multi-scale schedules:} Different schedules for different scales
\end{itemize}

\subsection{Better Training Techniques}

\begin{itemize}
    \item \textbf{Curriculum learning:} Start with easy examples
    \item \textbf{Data augmentation:} Use data augmentation during training
    \item \textbf{Regularization:} Better regularization techniques
\end{itemize}

\section{Applications}

\subsection{Text Generation}

\begin{itemize}
    \item \textbf{Story generation:} Generate coherent stories
    \item \textbf{Poetry generation:} Generate creative poetry
    \item \textbf{Code generation:} Generate programming code
\end{itemize}

\subsection{Image Generation}

\begin{itemize}
    \item \textbf{High-resolution images:} Generate high-quality images
    \item \textbf{Style transfer:} Transfer artistic styles
    \item \textbf{Image editing:} Edit existing images
\end{itemize}

\subsection{Scientific Applications}

\begin{itemize}
    \item \textbf{Molecular generation:} Generate new molecules
    \item \textbf{Protein design:} Design new proteins
    \item \textbf{Material discovery:} Discover new materials
\end{itemize}

\section{Future Directions}

\subsection{Theoretical Advances}

\begin{itemize}
    \item \textbf{Better understanding:} Theoretical analysis of discrete diffusion
    \item \textbf{Convergence guarantees:} Better convergence guarantees
    \item \textbf{Optimal schedules:} Find optimal noise schedules
\end{itemize}

\subsection{Architectural Improvements}

\begin{itemize}
    \item \textbf{More efficient models:} Reduce computational requirements
    \item \textbf{Better conditioning:} Improved conditional generation
    \item \textbf{Multi-modal models:} Handle multiple modalities
\end{itemize}

\subsection{Applications}

\begin{itemize}
    \item \textbf{Longer sequences:} Handle longer text sequences
    \item \textbf{Better quality:} Improve generation quality
    \item \textbf{New domains:} Apply to new domains and tasks
\end{itemize}

\section{Conclusion}

Discrete diffusion models extend the powerful diffusion framework to discrete data:

\textbf{Key contributions:}
\begin{itemize}
    \item \textbf{Discrete extension:} Extend diffusion to discrete state spaces
    \item \textbf{Flexible framework:} Handle various types of discrete data
    \item \textbf{Exact likelihood:} Can compute exact likelihood
    \item \textbf{Better control:} More control over generation process
\end{itemize}

\textbf{Key insights:}
\begin{itemize}
    \item \textbf{Transition matrices:} Define noise through transition matrices
    \item \textbf{Training objectives:} Use cross-entropy loss for training
    \item \textbf{Generation process:} Iterative denoising for generation
    \item \textbf{Applications:} Wide range of applications to discrete data
\end{itemize}

\textbf{Future work:}
\begin{itemize}
    \item \textbf{Efficiency:} More efficient training and inference
    \item \textbf{Quality:} Better generation quality
    \item \textbf{Theoretical understanding:} Better theoretical understanding
    \item \textbf{Applications:} New applications and domains
\end{itemize}

Discrete diffusion models provide a powerful framework for generative modeling of discrete data, with ongoing research exploring new applications and improvements.

\chapter{World Models: Learning Internal Representations of Environments}

\section{Introduction}

World models represent a paradigm shift in reinforcement learning and generative modeling, where the goal is to learn an internal model of the environment that can be used for planning, imagination, and decision-making. Unlike traditional approaches that learn policies directly from environment interactions, world models learn to predict future states and rewards, enabling agents to plan and act in imagined scenarios.

This chapter provides a comprehensive treatment of world models, covering theoretical foundations, practical implementations, and applications to various domains.

\section{Mathematical Foundations}

\subsection{World Model Definition}

A world model is a learned representation of the environment that can predict future states and rewards given current states and actions:
\begin{equation}
    p(s_{t+1}, r_t | s_t, a_t, h_t)
\end{equation}

where:
\begin{itemize}
    \item $s_t$ is the state at time $t$
    \item $a_t$ is the action at time $t$
    \item $r_t$ is the reward at time $t$
    \item $h_t$ is the history up to time $t$
\end{itemize}

\subsection{Components of World Models}

A world model typically consists of three components:

\textbf{1. Encoder:} Maps observations to latent states
\begin{equation}
    z_t = \text{Encoder}(o_t)
\end{equation}

\textbf{2. Dynamics Model:} Predicts next latent state
\begin{equation}
    z_{t+1} = \text{Dynamics}(z_t, a_t)
\end{equation}

\textbf{3. Decoder:} Reconstructs observations from latent states
\begin{equation}
    \hat{o}_t = \text{Decoder}(z_t)
\end{equation}

\section{Training Objectives}

\subsection{Reconstruction Loss}

The encoder-decoder should accurately reconstruct observations:
\begin{equation}
    \mathcal{L}_{\text{recon}} = \mathbb{E}[\|o_t - \hat{o}_t\|^2]
\end{equation}

\subsection{Prediction Loss}

The dynamics model should accurately predict future states:
\begin{equation}
    \mathcal{L}_{\text{pred}} = \mathbb{E}[\|z_{t+1} - \hat{z}_{t+1}\|^2]
\end{equation}

\subsection{Reward Prediction}

The model should also predict rewards:
\begin{equation}
    \mathcal{L}_{\text{reward}} = \mathbb{E}[\|r_t - \hat{r}_t\|^2]
\end{equation}

\section{Architecture Designs}

\subsection{Recurrent World Models}

Use recurrent neural networks to model temporal dependencies:
\begin{align}
    h_t &= \text{RNN}(h_{t-1}, z_t, a_t) \\
    z_{t+1} &= \text{Linear}(h_t) \\
    \hat{o}_t &= \text{Decoder}(z_t)
\end{align}

\subsection{Transformer-based World Models}

Use transformers for long-range dependencies:
\begin{align}
    z_{t+1} &= \text{Transformer}(z_1, \ldots, z_t, a_1, \ldots, a_t) \\
    \hat{o}_t &= \text{Decoder}(z_t)
\end{align}

\subsection{VAE-based World Models}

Use variational autoencoders for latent representations:
\begin{align}
    q_\phi(z_t | o_t) &= \mathcal{N}(\mu_\phi(o_t), \sigma_\phi^2(o_t)) \\
    p_\theta(o_t | z_t) &= \mathcal{N}(\mu_\theta(z_t), \sigma_\theta^2(z_t))
\end{align}

\section{Planning with World Models}

\subsection{Model Predictive Control (MPC)}

Use the world model for planning:
\begin{enumerate}
    \item Start with current state $s_0$
    \item For each action sequence $(a_1, \ldots, a_H)$:
    \begin{itemize}
        \item Roll out the world model: $s_{t+1} = f(s_t, a_t)$
        \item Compute cumulative reward: $R = \sum_{t=1}^H r_t$
    \end{itemize}
    \item Choose action sequence with highest reward
\end{enumerate}

\subsection{Tree Search}

Use tree search algorithms with the world model:
\begin{itemize}
    \item \textbf{MCTS:} Monte Carlo Tree Search
    \item \textbf{AlphaZero:} Combine MCTS with neural networks
    \item \textbf{Imagination:} Use world model for imagination
\end{itemize}

\section{Applications}

\subsection{Atari Games}

World models have been successfully applied to Atari games:
\begin{itemize}
    \item \textbf{World Models:} Learn compact representations of game states
    \item \textbf{Planning:} Use world model for planning
    \item \textbf{Imagination:} Generate imagined game scenarios
\end{itemize}

\subsection{Robotics}

World models are particularly useful in robotics:
\begin{itemize}
    \item \textbf{Sim-to-real:} Transfer from simulation to real world
    \item \textbf{Planning:} Plan robot actions using world model
    \item \textbf{Imagination:} Imagine consequences of actions
\end{itemize}

\subsection{Autonomous Driving}

World models can help with autonomous driving:
\begin{itemize}
    \item \textbf{Prediction:} Predict other vehicles' behavior
    \item \textbf{Planning:} Plan safe driving trajectories
    \item \textbf{Simulation:} Generate realistic driving scenarios
\end{itemize}

\section{Challenges and Limitations}

\subsection{Model Accuracy}

\begin{itemize}
    \item \textbf{Compounding errors:} Errors accumulate over time
    \item \textbf{Distribution shift:} Model may not generalize to new situations
    \item \textbf{Uncertainty:} Model may be overconfident in predictions
\end{itemize}

\subsection{Computational Challenges}

\begin{itemize}
    \item \textbf{Training time:} Can be expensive to train
    \item \textbf{Memory requirements:} Need to store large models
    \item \textbf{Inference speed:} May be slow for real-time applications
\end{itemize}

\section{Recent Advances}

\subsection{Improved Architectures}

\begin{itemize}
    \item \textbf{Transformer-based:} Use transformers for better long-range dependencies
    \item \textbf{Attention mechanisms:} Use attention to focus on relevant information
    \item \textbf{Multi-scale models:} Handle different temporal scales
\end{itemize}

\subsection{Better Training Techniques}

\begin{itemize}
    \item \textbf{Curriculum learning:} Start with easy tasks
    \item \textbf{Data augmentation:} Use data augmentation during training
    \item \textbf{Regularization:} Better regularization techniques
\end{itemize}

\section{Future Directions}

\subsection{Theoretical Advances}

\begin{itemize}
    \item \textbf{Better understanding:} Theoretical analysis of world models
    \item \textbf{Convergence guarantees:} Better convergence guarantees
    \item \textbf{Generalization:} Better understanding of generalization
\end{itemize}

\subsection{Architectural Improvements}

\begin{itemize}
    \item \textbf{More efficient models:} Reduce computational requirements
    \item \textbf{Better representations:} Learn better latent representations
    \item \textbf{Multi-modal models:} Handle multiple modalities
\end{itemize}

\subsection{Applications}

\begin{itemize}
    \item \textbf{New domains:} Apply to new domains and tasks
    \item \textbf{Better planning:} Improve planning algorithms
    \item \textbf{Real-world deployment:} Deploy in real-world applications
\end{itemize}

\section{Conclusion}

World models represent a powerful approach to learning internal representations of environments:

\textbf{Key contributions:}
\begin{itemize}
    \item \textbf{Internal models:} Learn internal representations of environments
    \item \textbf{Planning:} Use models for planning and decision-making
    \item \textbf{Imagination:} Generate imagined scenarios
    \item \textbf{Transfer:} Transfer knowledge across tasks
\end{itemize}

\textbf{Key insights:}
\begin{itemize}
    \item \textbf{Representation learning:} Learn good representations of states
    \item \textbf{Dynamics modeling:} Model environment dynamics
    \item \textbf{Planning:} Use models for planning
    \item \textbf{Imagination:} Use models for imagination
\end{itemize}

\textbf{Future work:}
\begin{itemize}
    \item \textbf{Efficiency:} More efficient training and inference
    \item \textbf{Accuracy:} Better model accuracy
    \item \textbf{Theoretical understanding:} Better theoretical understanding
    \item \textbf{Applications:} New applications and domains
\end{itemize}

World models provide a powerful framework for learning internal representations of environments, with ongoing research exploring new applications and improvements.

\chapter{Reinforcement Learning: Learning to Act in Dynamic Environments}

\section{Introduction}

Reinforcement Learning (RL) is a paradigm of machine learning where an agent learns to make decisions by interacting with an environment. Unlike supervised learning, RL does not require labeled examples of correct behavior. Instead, the agent learns through trial and error, receiving feedback in the form of rewards or penalties.

This chapter provides a comprehensive treatment of reinforcement learning, covering theoretical foundations, practical algorithms, and applications to various domains.

\section{Mathematical Foundations}

\subsection{Markov Decision Process (MDP)}

An MDP is defined by the tuple $(S, A, P, R, \gamma)$ where:
\begin{itemize}
    \item $S$ is the state space
    \item $A$ is the action space
    \item $P(s'|s,a)$ is the transition probability
    \item $R(s,a)$ is the reward function
    \item $\gamma \in [0,1]$ is the discount factor
\end{itemize}

\subsection{Policy and Value Functions}

\textbf{Policy:} A policy $\pi(a|s)$ defines the probability of taking action $a$ in state $s$.

\textbf{Value Function:} The value function $V^\pi(s)$ represents the expected cumulative reward starting from state $s$ and following policy $\pi$:
\begin{equation}
    V^\pi(s) = \mathbb{E}_\pi\left[\sum_{t=0}^\infty \gamma^t r_t | s_0 = s\right]
\end{equation}

\textbf{Q-Function:} The Q-function $Q^\pi(s,a)$ represents the expected cumulative reward starting from state $s$, taking action $a$, and then following policy $\pi$:
\begin{equation}
    Q^\pi(s,a) = \mathbb{E}_\pi\left[\sum_{t=0}^\infty \gamma^t r_t | s_0 = s, a_0 = a\right]
\end{equation}

\subsection{Bellman Equations}

The value function satisfies the Bellman equation:
\begin{equation}
    V^\pi(s) = \sum_a \pi(a|s) \sum_{s'} P(s'|s,a)[R(s,a) + \gamma V^\pi(s')]
\end{equation}

The Q-function satisfies:
\begin{equation}
    Q^\pi(s,a) = R(s,a) + \gamma \sum_{s'} P(s'|s,a) \sum_{a'} \pi(a'|s') Q^\pi(s',a')
\end{equation}

\section{Value-Based Methods}

\subsection{Value Iteration}

Value iteration updates the value function iteratively:
\begin{equation}
    V_{k+1}(s) = \max_a \sum_{s'} P(s'|s,a)[R(s,a) + \gamma V_k(s')]
\end{equation}

\subsection{Policy Iteration}

Policy iteration alternates between policy evaluation and policy improvement:
\begin{enumerate}
    \item \textbf{Policy Evaluation:} Compute $V^\pi$ for current policy $\pi$
    \item \textbf{Policy Improvement:} Update policy to be greedy with respect to $V^\pi$
\end{enumerate}

\subsection{Q-Learning}

Q-learning is a model-free algorithm that learns the Q-function:
\begin{equation}
    Q(s,a) \leftarrow Q(s,a) + \alpha[r + \gamma \max_{a'} Q(s',a') - Q(s,a)]
\end{equation}

\subsection{Deep Q-Networks (DQN)}

DQN uses neural networks to approximate the Q-function:
\begin{equation}
    Q(s,a) \approx Q_\theta(s,a)
\end{equation}

\textbf{Training objective:}
\begin{equation}
    \mathcal{L}(\theta) = \mathbb{E}[(r + \gamma \max_{a'} Q_{\theta^-}(s',a') - Q_\theta(s,a))^2]
\end{equation}

where $\theta^-$ is the target network.

\section{Policy-Based Methods}

\subsection{Policy Gradient}

The policy gradient theorem states:
\begin{equation}
    \nabla_\theta J(\theta) = \mathbb{E}_\pi\left[\sum_{t=0}^T \nabla_\theta \log \pi(a_t|s_t) A_t\right]
\end{equation}

where $A_t$ is the advantage function.

\subsection{REINFORCE}

REINFORCE is a simple policy gradient algorithm:
\begin{equation}
    \nabla_\theta J(\theta) = \mathbb{E}_\pi\left[\sum_{t=0}^T \nabla_\theta \log \pi(a_t|s_t) R_t\right]
\end{equation}

\subsection{Actor-Critic Methods}

Actor-critic methods combine policy and value function learning:
\begin{align}
    \nabla_\theta J(\theta) &= \mathbb{E}_\pi\left[\sum_{t=0}^T \nabla_\theta \log \pi(a_t|s_t) A_t\right] \\
    A_t &= Q(s_t,a_t) - V(s_t)
\end{align}

\section{Advanced Algorithms}

\subsection{Proximal Policy Optimization (PPO)}

PPO is a policy gradient algorithm that uses clipping to prevent large policy updates:
\begin{equation}
    \mathcal{L}(\theta) = \mathbb{E}\left[\min(r_t(\theta) A_t, \text{clip}(r_t(\theta), 1-\epsilon, 1+\epsilon) A_t)\right]
\end{equation}

where $r_t(\theta) = \frac{\pi_\theta(a_t|s_t)}{\pi_{\theta_{old}}(a_t|s_t)}$.

\subsection{Trust Region Policy Optimization (TRPO)}

TRPO uses a trust region to ensure policy updates don't degrade performance:
\begin{equation}
    \max_\theta \mathbb{E}\left[\frac{\pi_\theta(a|s)}{\pi_{\theta_{old}}(a|s)} A_t\right]
\end{equation}

subject to:
\begin{equation}
    \mathbb{E}[D_{KL}(\pi_{\theta_{old}}(\cdot|s) \| \pi_\theta(\cdot|s))] \leq \delta
\end{equation}

\subsection{Soft Actor-Critic (SAC)}

SAC is an off-policy algorithm that maximizes entropy:
\begin{equation}
    J(\pi) = \mathbb{E}\left[\sum_{t=0}^T \gamma^t (r_t + \alpha \mathcal{H}(\pi(\cdot|s_t)))\right]
\end{equation}

\section{Function Approximation}

\subsection{Neural Networks}

Neural networks can approximate value functions and policies:
\begin{align}
    V(s) &\approx V_\theta(s) \\
    Q(s,a) &\approx Q_\theta(s,a) \\
    \pi(a|s) &\approx \pi_\theta(a|s)
\end{align}

\subsection{Experience Replay}

Experience replay stores past experiences and samples them randomly:
\begin{itemize}
    \item \textbf{Buffer:} Store $(s_t, a_t, r_t, s_{t+1})$ tuples
    \item \textbf{Sampling:} Sample random batches for training
    \item \textbf{Benefits:} Reduces correlation, improves sample efficiency
\end{itemize}

\subsection{Target Networks}

Target networks provide stable targets for training:
\begin{itemize}
    \item \textbf{Main network:} Updated every step
    \item \textbf{Target network:} Updated less frequently
    \item \textbf{Benefits:} Reduces instability, improves convergence
\end{itemize}

\section{Exploration Strategies}

\subsection{Epsilon-Greedy}

$\epsilon$-greedy exploration balances exploration and exploitation:
\begin{equation}
    a_t = \begin{cases}
        \arg\max_a Q(s_t,a) & \text{with probability } 1-\epsilon \\
        \text{random action} & \text{with probability } \epsilon
    \end{cases}
\end{equation}

\subsection{Upper Confidence Bound (UCB)}

UCB balances exploration and exploitation based on uncertainty:
\begin{equation}
    a_t = \arg\max_a \left[Q(s_t,a) + c\sqrt{\frac{\log t}{N_t(a)}}\right]
\end{equation}

\subsection{Thompson Sampling}

Thompson sampling uses Bayesian inference for exploration:
\begin{equation}
    a_t = \arg\max_a Q(s_t,a) + \sigma(s_t,a) \cdot \epsilon
\end{equation}

where $\epsilon \sim \mathcal{N}(0,1)$.

\section{Applications}

\subsection{Games}

RL has been successfully applied to games:
\begin{itemize}
    \item \textbf{Atari:} DQN achieved human-level performance
    \item \textbf{Go:} AlphaGo defeated world champion
    \item \textbf{Chess:} AlphaZero achieved superhuman performance
\end{itemize}

\subsection{Robotics}

RL is widely used in robotics:
\begin{itemize}
    \item \textbf{Manipulation:} Learn to manipulate objects
    \item \textbf{Locomotion:} Learn to walk and run
    \item \textbf{Navigation:} Learn to navigate environments
\end{itemize}

\subsection{Autonomous Systems}

RL is used in autonomous systems:
\begin{itemize}
    \item \textbf{Autonomous driving:} Learn to drive safely
    \item \textbf{Drone control:} Learn to fly drones
    \item \textbf{Resource allocation:} Learn to allocate resources
\end{itemize}

\section{Challenges and Limitations}

\subsection{Sample Efficiency}

\begin{itemize}
    \item \textbf{High sample complexity:} RL often requires many samples
    \item \textbf{Exploration:} Need to explore environment effectively
    \item \textbf{Transfer learning:} Difficult to transfer knowledge across tasks
\end{itemize}

\subsection{Stability}

\begin{itemize}
    \item \textbf{Training instability:} RL training can be unstable
    \item \textbf{Hyperparameter sensitivity:} Performance sensitive to hyperparameters
    \item \textbf{Convergence:} No guarantee of convergence to optimal policy
\end{itemize}

\subsection{Safety}

\begin{itemize}
    \item \textbf{Safe exploration:} Need to explore safely
    \item \textbf{Robustness:} Policies may not be robust to distribution shift
    \item \textbf{Interpretability:} Difficult to interpret learned policies
\end{itemize}

\section{Recent Advances}

\subsection{Model-Based RL}

Model-based RL learns a model of the environment:
\begin{itemize}
    \item \textbf{World models:} Learn internal models of environments
    \item \textbf{Planning:} Use models for planning
    \item \textbf{Sample efficiency:} Can be more sample efficient
\end{itemize}

\subsection{Multi-Agent RL}

Multi-agent RL considers multiple agents:
\begin{itemize}
    \item \textbf{Cooperation:} Agents can cooperate
    \item \textbf{Competition:} Agents can compete
    \item \textbf{Coordination:} Agents need to coordinate
\end{itemize}

\subsection{Hierarchical RL}

Hierarchical RL uses multiple levels of abstraction:
\begin{itemize}
    \item \textbf{High-level policies:} Abstract actions
    \item \textbf{Low-level policies:} Concrete actions
    \item \textbf{Decomposition:} Break down complex tasks
\end{itemize}

\section{Future Directions}

\subsection{Theoretical Advances}

\begin{itemize}
    \item \textbf{Better convergence guarantees:} Improve theoretical understanding
    \item \textbf{Sample complexity:} Reduce sample complexity
    \item \textbf{Generalization:} Better understanding of generalization
\end{itemize}

\subsection{Algorithmic Improvements}

\begin{itemize}
    \item \textbf{More efficient algorithms:} Reduce computational requirements
    \item \textbf{Better exploration:} Improve exploration strategies
    \item \textbf{Transfer learning:} Better transfer across tasks
\end{itemize}

\subsection{Applications}

\begin{itemize}
    \item \textbf{Real-world deployment:} Deploy in real-world applications
    \item \textbf{New domains:} Apply to new domains and tasks
    \item \textbf{Safety:} Ensure safe operation
\end{itemize}

\section{Conclusion}

Reinforcement learning provides a powerful framework for learning to act in dynamic environments:

\textbf{Key contributions:}
\begin{itemize}
    \item \textbf{Learning from interaction:} Learn from trial and error
    \item \textbf{Decision making:} Make decisions in dynamic environments
    \item \textbf{Optimization:} Optimize long-term rewards
    \item \textbf{Autonomy:} Enable autonomous systems
\end{itemize}

\textbf{Key insights:}
\begin{itemize}
    \item \textbf{Value functions:} Represent expected future rewards
    \item \textbf{Policy optimization:} Learn to improve policies
    \item \textbf{Exploration:} Balance exploration and exploitation
    \item \textbf{Function approximation:} Use neural networks for complex environments
\end{itemize}

\textbf{Future work:}
\begin{itemize}
    \item \textbf{Efficiency:} More sample-efficient algorithms
    \item \textbf{Safety:} Ensure safe operation
    \item \textbf{Theoretical understanding:} Better theoretical understanding
    \item \textbf{Applications:} New applications and domains
\end{itemize}

Reinforcement learning continues to drive innovation in artificial intelligence, with ongoing research exploring new algorithms, applications, and theoretical foundations.

\chapter{Reasoning: From Symbolic Logic to Neural Reasoning}

\section{Introduction}

Reasoning represents the ability to draw logical conclusions from given premises, a fundamental aspect of human intelligence that has been a long-standing challenge in artificial intelligence. While traditional symbolic approaches to reasoning have been successful in limited domains, recent advances in neural networks have opened new possibilities for learning to reason from data.

This chapter provides a comprehensive treatment of reasoning in AI, covering both symbolic and neural approaches, their mathematical foundations, and applications to various domains.

\section{Symbolic Reasoning}

\subsection{Propositional Logic}

Propositional logic deals with propositions that can be true or false:
\begin{itemize}
    \item \textbf{Atomic propositions:} $p, q, r, \ldots$
    \item \textbf{Logical connectives:} $\neg, \land, \lor, \rightarrow, \leftrightarrow$
    \item \textbf{Truth tables:} Define truth values for all combinations
\end{itemize}

\subsection{Predicate Logic}

Predicate logic extends propositional logic with quantifiers:
\begin{itemize}
    \item \textbf{Predicates:} $P(x), Q(x,y), \ldots$
    \item \textbf{Quantifiers:} $\forall$ (for all), $\exists$ (there exists)
    \item \textbf{Variables:} $x, y, z, \ldots$
\end{itemize}

\subsection{Inference Rules}

\textbf{Modus Ponens:}
\begin{equation}
    \frac{P \rightarrow Q, P}{Q}
\end{equation}

\textbf{Modus Tollens:}
\begin{equation}
    \frac{P \rightarrow Q, \neg Q}{\neg P}
\end{equation}

\textbf{Resolution:}
\begin{equation}
    \frac{P \lor Q, \neg P \lor R}{Q \lor R}
\end{equation}

\section{Neural Reasoning}

\subsection{Neural Symbolic Reasoning}

Neural symbolic reasoning combines neural networks with symbolic reasoning:
\begin{itemize}
    \item \textbf{Neural encoding:} Encode symbolic knowledge in neural networks
    \item \textbf{Symbolic reasoning:} Perform logical inference
    \item \textbf{Neural decoding:} Convert results back to symbolic form
\end{itemize}

\subsection{Graph Neural Networks for Reasoning}

Graph Neural Networks (GNNs) can reason about relational data:
\begin{equation}
    h_v^{(l+1)} = \text{AGGREGATE}(\{h_u^{(l)} : u \in \mathcal{N}(v)\})
\end{equation}

where $\mathcal{N}(v)$ is the neighborhood of node $v$.

\subsection{Attention Mechanisms for Reasoning}

Attention mechanisms can focus on relevant information for reasoning:
\begin{equation}
    \text{Attention}(Q, K, V) = \text{softmax}\left(\frac{QK^T}{\sqrt{d_k}}\right)V
\end{equation}

\section{Mathematical Reasoning}

\subsection{Arithmetic Reasoning}

Neural networks can learn to perform arithmetic operations:
\begin{itemize}
    \item \textbf{Addition:} $a + b = c$
    \item \textbf{Subtraction:} $a - b = c$
    \item \textbf{Multiplication:} $a \times b = c$
    \item \textbf{Division:} $a \div b = c$
\end{itemize}

\subsection{Algebraic Reasoning}

Neural networks can solve algebraic equations:
\begin{equation}
    ax + b = c \Rightarrow x = \frac{c - b}{a}
\end{equation}

\subsection{Geometric Reasoning}

Neural networks can reason about geometric properties:
\begin{itemize}
    \item \textbf{Shapes:} Identify and classify shapes
    \item \textbf{Relations:} Understand spatial relationships
    \item \textbf{Transformations:} Apply geometric transformations
\end{itemize}

\section{Logical Reasoning}

\subsection{Propositional Reasoning}

Neural networks can learn propositional logic:
\begin{itemize}
    \item \textbf{Truth tables:} Learn truth values for logical connectives
    \item \textbf{Inference:} Perform logical inference
    \item \textbf{Proofs:} Generate logical proofs
\end{itemize}

\subsection{Predicate Reasoning}

Neural networks can handle predicate logic:
\begin{itemize}
    \item \textbf{Quantifiers:} Handle universal and existential quantifiers
    \item \textbf{Unification:} Find substitutions for variables
    \item \textbf{Resolution:} Perform resolution theorem proving
\end{itemize}

\section{Common Sense Reasoning}

\subsection{Physical Reasoning}

Neural networks can reason about physical properties:
\begin{itemize}
    \item \textbf{Gravity:} Understand effects of gravity
    \item \textbf{Momentum:} Understand conservation of momentum
    \item \textbf{Collisions:} Predict outcomes of collisions
\end{itemize}

\subsection{Causal Reasoning}

Neural networks can learn causal relationships:
\begin{itemize}
    \item \textbf{Causality:} Identify cause-effect relationships
    \item \textbf{Intervention:} Predict effects of interventions
    \item \textbf{Counterfactuals:} Reason about counterfactual scenarios
\end{itemize}

\section{Reasoning in Language}

\subsection{Natural Language Inference}

Neural networks can perform natural language inference:
\begin{itemize}
    \item \textbf{Entailment:} Determine if one sentence entails another
    \item \textbf{Contradiction:} Identify contradictory statements
    \item \textbf{Neutral:} Determine if statements are unrelated
\end{itemize}

\subsection{Question Answering}

Neural networks can answer questions by reasoning:
\begin{itemize}
    \item \textbf{Reading comprehension:} Answer questions about text
    \item \textbf{Multi-hop reasoning:} Answer questions requiring multiple reasoning steps
    \item \textbf{Commonsense QA:} Answer questions requiring commonsense knowledge
\end{itemize}

\section{Reasoning in Vision}

\subsection{Spatial Reasoning}

Neural networks can reason about spatial relationships:
\begin{itemize}
    \item \textbf{Position:} Understand object positions
    \item \textbf{Orientation:} Understand object orientations
    \item \textbf{Relations:} Understand spatial relationships between objects
\end{itemize}

\subsection{Temporal Reasoning}

Neural networks can reason about temporal sequences:
\begin{itemize}
    \item \textbf{Order:} Understand temporal order of events
    \item \textbf{Duration:} Understand duration of events
    \item \textbf{Causality:} Understand temporal causality
\end{itemize}

\section{Reasoning Architectures}

\subsection{Memory Networks}

Memory networks use external memory for reasoning:
\begin{itemize}
    \item \textbf{Memory:} Store facts and information
    \item \textbf{Attention:} Attend to relevant memories
    \item \textbf{Reasoning:} Perform reasoning over memories
\end{itemize}

\subsection{Neural Turing Machines}

Neural Turing Machines combine neural networks with external memory:
\begin{itemize}
    \item \textbf{Controller:} Neural network controller
    \item \textbf{Memory:} External memory matrix
    \item \textbf{Heads:} Read and write heads
\end{itemize}

\subsection{Transformer-based Reasoning}

Transformers can perform reasoning through attention:
\begin{itemize}
    \item \textbf{Self-attention:} Attend to different parts of input
    \item \textbf{Cross-attention:} Attend to different modalities
    \item \textbf{Multi-head attention:} Multiple attention mechanisms
\end{itemize}

\section{Evaluation Metrics}

\subsection{Accuracy}

Standard accuracy measures:
\begin{itemize}
    \item \textbf{Exact match:} Exact match with ground truth
    \item \textbf{F1 score:} Harmonic mean of precision and recall
    \item \textbf{BLEU score:} N-gram overlap for text generation
\end{itemize}

\subsection{Reasoning-specific Metrics}

\begin{itemize}
    \item \textbf{Logical consistency:} Check logical consistency of outputs
    \item \textbf{Proof correctness:} Verify correctness of proofs
    \item \textbf{Step-by-step accuracy:} Accuracy of individual reasoning steps
\end{itemize}

\section{Challenges and Limitations}

\subsection{Symbolic vs. Neural}

\begin{itemize}
    \item \textbf{Symbolic:} Interpretable but limited scalability
    \item \textbf{Neural:} Scalable but less interpretable
    \item \textbf{Hybrid:} Combine benefits of both approaches
\end{itemize}

\subsection{Generalization}

\begin{itemize}
    \item \textbf{Domain transfer:} Transfer reasoning across domains
    \item \textbf{Compositionality:} Compose simple reasoning steps
    \item \textbf{Abstraction:} Reason at different levels of abstraction
\end{itemize}

\subsection{Interpretability}

\begin{itemize}
    \item \textbf{Black box:} Neural networks are often black boxes
    \item \textbf{Explanation:} Need to explain reasoning steps
    \item \textbf{Debugging:} Difficult to debug reasoning errors
\end{itemize}

\section{Recent Advances}

\subsection{Large Language Models}

Large language models have shown impressive reasoning capabilities:
\begin{itemize}
    \item \textbf{GPT:} Generative pre-trained transformer
    \item \textbf{BERT:} Bidirectional encoder representations
    \item \textbf{T5:} Text-to-text transfer transformer
\end{itemize}

\subsection{Chain-of-Thought Reasoning}

Chain-of-thought reasoning generates step-by-step reasoning:
\begin{itemize}
    \item \textbf{Step-by-step:} Generate reasoning steps
    \item \textbf{Verification:} Verify each reasoning step
    \item \textbf{Correction:} Correct errors in reasoning
\end{itemize}

\subsection{Program Synthesis}

Program synthesis generates programs to solve reasoning tasks:
\begin{itemize}
    \item \textbf{Code generation:} Generate code for reasoning tasks
    \item \textbf{Verification:} Verify correctness of generated code
    \item \textbf{Optimization:} Optimize generated programs
\end{itemize}

\section{Applications}

\subsection{Scientific Discovery}

\begin{itemize}
    \item \textbf{Hypothesis generation:} Generate scientific hypotheses
    \item \textbf{Experiment design:} Design scientific experiments
    \item \textbf{Data analysis:} Analyze scientific data
\end{itemize}

\subsection{Legal Reasoning}

\begin{itemize}
    \item \textbf{Case analysis:} Analyze legal cases
    \item \textbf{Precedent finding:} Find relevant legal precedents
    \item \textbf{Argument construction:} Construct legal arguments
\end{itemize}

\subsection{Medical Diagnosis}

\begin{itemize}
    \item \textbf{Symptom analysis:} Analyze patient symptoms
    \item \textbf{Differential diagnosis:} Generate differential diagnoses
    \item \textbf{Treatment planning:} Plan treatments
\end{itemize}

\section{Future Directions}

\subsection{Theoretical Advances}

\begin{itemize}
    \item \textbf{Formal verification:} Verify correctness of reasoning systems
    \item \textbf{Theoretical guarantees:} Provide theoretical guarantees
    \item \textbf{Convergence:} Understand convergence properties
\end{itemize}

\subsection{Architectural Improvements}

\begin{itemize}
    \item \textbf{Better architectures:} Design better reasoning architectures
    \item \textbf{Efficiency:} Improve computational efficiency
    \item \textbf{Scalability:} Scale to larger problems
\end{itemize}

\subsection{Applications}

\begin{itemize}
    \item \textbf{New domains:} Apply to new domains and tasks
    \item \textbf{Real-world deployment:} Deploy in real-world applications
    \item \textbf{Human-AI collaboration:} Collaborate with humans
\end{itemize}

\section{Conclusion}

Reasoning represents a fundamental aspect of intelligence that has seen significant advances in recent years:

\textbf{Key contributions:}
\begin{itemize}
    \item \textbf{Symbolic reasoning:} Traditional logical reasoning approaches
    \item \textbf{Neural reasoning:} Neural network-based reasoning
    \item \textbf{Hybrid approaches:} Combining symbolic and neural methods
    \item \textbf{Applications:} Wide range of applications across domains
\end{itemize}

\textbf{Key insights:}
\begin{itemize}
    \item \textbf{Representation:} Good representations are crucial for reasoning
    \item \textbf{Attention:} Attention mechanisms enable focused reasoning
    \item \textbf{Memory:} External memory can enhance reasoning capabilities
    \item \textbf{Compositionality:} Complex reasoning can be composed from simple steps
\end{itemize}

\textbf{Future work:}
\begin{itemize}
    \item \textbf{Interpretability:} Make reasoning more interpretable
    \item \textbf{Generalization:} Improve generalization across domains
    \item \textbf{Efficiency:} More efficient reasoning algorithms
    \item \textbf{Applications:} New applications and domains
\end{itemize}

Reasoning continues to be a central challenge in AI, with ongoing research exploring new approaches, architectures, and applications.

\chapter{Intelligent Agents: From Reactive to Autonomous Systems}

\section{Introduction}

Intelligent agents represent autonomous entities that can perceive their environment, reason about it, and take actions to achieve their goals. The concept of agents has evolved from simple reactive systems to sophisticated autonomous entities capable of complex reasoning, learning, and decision-making.

This chapter provides a comprehensive treatment of intelligent agents, covering their theoretical foundations, architectures, learning mechanisms, and applications to various domains.

\section{Agent Architecture}

\subsection{Basic Agent Model}

An agent can be characterized by its:
\begin{itemize}
    \item \textbf{Perception:} How it perceives the environment
    \item \textbf{Reasoning:} How it reasons about the environment
    \item \textbf{Action:} How it acts in the environment
    \item \textbf{Goals:} What it aims to achieve
\end{itemize}

\subsection{Agent-Environment Interaction}

The agent-environment interaction can be modeled as:
\begin{equation}
    \text{Agent: } s_t \rightarrow a_t
\end{equation}
\begin{equation}
    \text{Environment: } (s_t, a_t) \rightarrow (s_{t+1}, r_t)
\end{equation}

where $s_t$ is the state, $a_t$ is the action, and $r_t$ is the reward.

\section{Types of Agents}

\subsection{Reactive Agents}

Reactive agents respond directly to environmental stimuli:
\begin{equation}
    a_t = \pi(s_t)
\end{equation}

where $\pi$ is a policy function.

\subsection{Deliberative Agents}

Deliberative agents plan before acting:
\begin{enumerate}
    \item \textbf{Goal formulation:} Determine what to achieve
    \item \textbf{Planning:} Generate a plan to achieve the goal
    \item \textbf{Execution:} Execute the plan
\end{enumerate}

\subsection{Hybrid Agents}

Hybrid agents combine reactive and deliberative capabilities:
\begin{itemize}
    \item \textbf{Reactive layer:} Handle immediate responses
    \item \textbf{Deliberative layer:} Handle complex planning
    \item \textbf{Coordination:} Coordinate between layers
\end{itemize}

\section{Learning in Agents}

\subsection{Reinforcement Learning}

Agents can learn through reinforcement learning:
\begin{equation}
    Q(s,a) \leftarrow Q(s,a) + \alpha[r + \gamma \max_{a'} Q(s',a') - Q(s,a)]
\end{equation}

\subsection{Imitation Learning}

Agents can learn by imitating expert behavior:
\begin{equation}
    \pi^* = \arg\min_\pi \mathbb{E}_{s \sim d^\pi}[\ell(\pi(s), \pi^*(s))]
\end{equation}

\subsection{Transfer Learning}

Agents can transfer knowledge across tasks:
\begin{itemize}
    \item \textbf{Feature transfer:} Transfer learned features
    \item \textbf{Policy transfer:} Transfer learned policies
    \item \textbf{Model transfer:} Transfer learned models
\end{itemize}

\section{Multi-Agent Systems}

\subsection{Cooperative Agents}

Cooperative agents work together to achieve common goals:
\begin{itemize}
    \item \textbf{Coordination:} Coordinate actions among agents
    \item \textbf{Communication:} Communicate information
    \item \textbf{Consensus:} Reach consensus on decisions
\end{itemize}

\subsection{Competitive Agents}

Competitive agents compete against each other:
\begin{itemize}
    \item \textbf{Game theory:} Use game-theoretic approaches
    \item \textbf{Nash equilibrium:} Find Nash equilibria
    \item \textbf{Adversarial training:} Train against opponents
\end{itemize}

\subsection{Mixed Environments}

Mixed environments contain both cooperative and competitive agents:
\begin{itemize}
    \item \textbf{Coalition formation:} Form coalitions
    \item \textbf{Negotiation:} Negotiate agreements
    \item \textbf{Conflict resolution:} Resolve conflicts
\end{itemize}

\section{Agent Communication}

\subsection{Message Passing}

Agents can communicate through message passing:
\begin{itemize}
    \item \textbf{Synchronous:} Messages sent and received simultaneously
    \item \textbf{Asynchronous:} Messages sent and received at different times
    \item \textbf{Reliable:} Guarantee message delivery
    \item \textbf{Unreliable:} No guarantee of message delivery
\end{itemize}

\subsection{Shared Memory}

Agents can share information through shared memory:
\begin{itemize}
    \item \textbf{Blackboard:} Shared blackboard for information
    \item \textbf{Databases:} Shared databases
    \item \textbf{Knowledge bases:} Shared knowledge bases
\end{itemize}

\subsection{Protocols}

Communication protocols define how agents interact:
\begin{itemize}
    \item \textbf{Request-Response:} Request information and receive response
    \item \textbf{Publish-Subscribe:} Publish information and subscribe to updates
    \item \textbf{Auction:} Bid for resources or tasks
\end{itemize}

\section{Agent Planning}

\subsection{Classical Planning}

Classical planning assumes:
\begin{itemize}
    \item \textbf{Deterministic:} Actions have deterministic effects
    \item \textbf{Fully observable:} Complete information about environment
    \item \textbf{Static:} Environment doesn't change during planning
\end{itemize}

\subsection{Probabilistic Planning}

Probabilistic planning handles uncertainty:
\begin{itemize}
    \item \textbf{Stochastic actions:} Actions have probabilistic effects
    \item \textbf{Partial observability:} Incomplete information about environment
    \item \textbf{Dynamic:} Environment can change during execution
\end{itemize}

\subsection{Hierarchical Planning}

Hierarchical planning uses multiple levels of abstraction:
\begin{itemize}
    \item \textbf{High-level plans:} Abstract plans
    \item \textbf{Low-level plans:} Concrete plans
    \item \textbf{Decomposition:} Break down complex tasks
\end{itemize}

\section{Agent Learning}

\subsection{Online Learning}

Agents can learn while acting:
\begin{itemize}
    \item \textbf{Exploration:} Explore new actions
    \item \textbf{Exploitation:} Use learned knowledge
    \item \textbf{Balance:} Balance exploration and exploitation
\end{itemize}

\subsection{Offline Learning}

Agents can learn from historical data:
\begin{itemize}
    \item \textbf{Batch learning:} Learn from batches of data
    \item \textbf{Incremental learning:} Learn incrementally from new data
    \item \textbf{Transfer learning:} Transfer knowledge across tasks
\end{itemize}

\subsection{Meta-Learning}

Agents can learn how to learn:
\begin{itemize}
    \item \textbf{Learning to learn:} Learn learning algorithms
    \item \textbf{Few-shot learning:} Learn from few examples
    \item \textbf{Rapid adaptation:} Quickly adapt to new tasks
\end{itemize}

\section{Agent Architectures}

\subsection{Belief-Desire-Intention (BDI)}

BDI agents have:
\begin{itemize}
    \item \textbf{Beliefs:} What the agent believes about the world
    \item \textbf{Desires:} What the agent wants to achieve
    \item \textbf{Intentions:} What the agent plans to do
\end{itemize}

\subsection{Cognitive Architectures}

Cognitive architectures model human-like cognition:
\begin{itemize}
    \item \textbf{ACT-R:} Adaptive Control of Thought-Rational
    \item \textbf{SOAR:} State, Operator, And Result
    \item \textbf{CLARION:} Connectionist Learning with Adaptive Rule Induction
\end{itemize}

\subsection{Neural Architectures}

Neural architectures use neural networks:
\begin{itemize}
    \item \textbf{Deep Q-Networks:} Deep neural networks for Q-learning
    \item \textbf{Policy networks:} Neural networks for policy learning
    \item \textbf{Value networks:} Neural networks for value function approximation
\end{itemize}

\section{Agent Applications}

\subsection{Robotics}

Agents are widely used in robotics:
\begin{itemize}
    \item \textbf{Autonomous robots:} Robots that operate independently
    \item \textbf{Swarm robotics:} Multiple robots working together
    \item \textbf{Human-robot interaction:} Robots that interact with humans
\end{itemize}

\subsection{Autonomous Systems}

Agents are used in autonomous systems:
\begin{itemize}
    \item \textbf{Autonomous vehicles:} Self-driving cars
    \item \textbf{Unmanned aerial vehicles:} Drones
    \item \textbf{Autonomous spacecraft:} Space exploration missions
\end{itemize}

\subsection{Software Agents}

Agents are used in software systems:
\begin{itemize}
    \item \textbf{Web agents:} Agents that browse the web
    \item \textbf{Email agents:} Agents that manage email
    \item \textbf{Personal assistants:} Agents that help users
\end{itemize}

\section{Agent Evaluation}

\subsection{Performance Metrics}

Agents can be evaluated using various metrics:
\begin{itemize}
    \item \textbf{Task completion:} How well agents complete tasks
    \item \textbf{Efficiency:} How efficiently agents use resources
    \item \textbf{Reliability:} How reliably agents perform
    \item \textbf{Safety:} How safely agents operate
\end{itemize}

\subsection{Benchmarks}

Various benchmarks exist for evaluating agents:
\begin{itemize}
    \item \textbf{Atari:} Video game playing
    \item \textbf{MuJoCo:} Continuous control tasks
    \item \textbf{StarCraft:} Real-time strategy games
    \item \textbf{Dota 2:} Multi-player online battle arena
\end{itemize}

\section{Challenges and Limitations}

\subsection{Scalability}

\begin{itemize}
    \item \textbf{State space:} Large state spaces can be challenging
    \item \textbf{Action space:} Large action spaces can be challenging
    \item \textbf{Multi-agent:} Coordination among many agents
\end{itemize}

\subsection{Safety}

\begin{itemize}
    \item \textbf{Safe exploration:} Explore safely in dangerous environments
    \item \textbf{Robustness:} Handle unexpected situations
    \item \textbf{Verification:} Verify agent behavior
\end{itemize}

\subsection{Ethics}

\begin{itemize}
    \item \textbf{Fairness:} Ensure fair treatment of all agents
    \item \textbf{Privacy:} Protect agent privacy
    \item \textbf{Transparency:} Make agent decisions transparent
\end{itemize}

\section{Recent Advances}

\subsection{Large Language Models as Agents}

Large language models can act as agents:
\begin{itemize}
    \item \textbf{Text-based interaction:} Interact through text
    \item \textbf{Tool use:} Use external tools
    \item \textbf{Planning:} Plan and execute tasks
\end{itemize}

\subsection{Multi-Modal Agents}

Multi-modal agents can handle multiple modalities:
\begin{itemize}
    \item \textbf{Vision:} Process visual information
    \item \textbf{Language:} Process linguistic information
    \item \textbf{Audio:} Process auditory information
\end{itemize}

\subsection{Embodied Agents}

Embodied agents have physical bodies:
\begin{itemize}
    \item \textbf{Robotics:} Physical robots
    \item \textbf{Virtual agents:} Virtual characters
    \item \textbf{Simulation:} Simulated environments
\end{itemize}

\section{Future Directions}

\subsection{Theoretical Advances}

\begin{itemize}
    \item \textbf{Formal verification:} Verify agent behavior
    \item \textbf{Theoretical guarantees:} Provide theoretical guarantees
    \item \textbf{Convergence:} Understand convergence properties
\end{itemize}

\subsection{Architectural Improvements}

\begin{itemize}
    \item \textbf{Better architectures:} Design better agent architectures
    \item \textbf{Efficiency:} Improve computational efficiency
    \item \textbf{Scalability:} Scale to larger problems
\end{itemize}

\subsection{Applications}

\begin{itemize}
    \item \textbf{New domains:} Apply to new domains and tasks
    \item \textbf{Real-world deployment:} Deploy in real-world applications
    \item \textbf{Human-AI collaboration:} Collaborate with humans
\end{itemize}

\section{Conclusion}

Intelligent agents represent a fundamental paradigm in artificial intelligence:

\textbf{Key contributions:}
\begin{itemize}
    \item \textbf{Autonomy:} Autonomous operation in dynamic environments
    \item \textbf{Learning:} Learn from experience and adapt
    \item \textbf{Reasoning:} Reason about environment and goals
    \item \textbf{Interaction:} Interact with environment and other agents
\end{itemize}

\textbf{Key insights:}
\begin{itemize}
    \item \textbf{Architecture:} Good agent architecture is crucial
    \item \textbf{Learning:} Learning enables adaptation and improvement
    \item \textbf{Communication:} Communication enables coordination
    \item \textbf{Planning:} Planning enables goal achievement
\end{itemize}

\textbf{Future work:}
\begin{itemize}
    \item \textbf{Safety:} Ensure safe operation
    \item \textbf{Efficiency:} More efficient agents
    \item \textbf{Theoretical understanding:} Better theoretical understanding
    \item \textbf{Applications:} New applications and domains
\end{itemize}

Intelligent agents continue to be a central focus in AI research, with ongoing work exploring new architectures, learning mechanisms, and applications.

\chapter{AI for Science: Accelerating Scientific Discovery}

\section{Introduction}

Artificial Intelligence has emerged as a powerful tool for accelerating scientific discovery across diverse domains, from physics and chemistry to biology and medicine. The integration of AI with scientific research has led to breakthroughs in drug discovery, materials science, climate modeling, and fundamental physics research.

This chapter provides a comprehensive treatment of AI for science, covering theoretical foundations, practical applications, and the transformative impact on scientific research.

\section{Mathematical Foundations}

\subsection{Scientific Data Representation}

Scientific data can be represented in various forms:
\begin{itemize}
    \item \textbf{Time series:} $x(t) = \{x_1, x_2, \ldots, x_T\}$
    \item \textbf{Images:} $I(x,y) \in \mathbb{R}^{H \times W \times C}$
    \item \textbf{Graphs:} $G = (V, E)$ where $V$ are nodes and $E$ are edges
    \item \textbf{Sequences:} $s = \{s_1, s_2, \ldots, s_n\}$ (e.g., DNA sequences)
\end{itemize}

\subsection{Physical Constraints}

AI models for science must respect physical constraints:
\begin{itemize}
    \item \textbf{Conservation laws:} Energy, momentum, mass conservation
    \item \textbf{Symmetries:} Rotational, translational, gauge symmetries
    \item \textbf{Causality:} Cause-effect relationships
    \item \textbf{Uncertainty:} Quantify and propagate uncertainty
\end{itemize}

\section{Drug Discovery}

\subsection{Molecular Representation}

Molecules can be represented in various ways:
\begin{itemize}
    \item \textbf{SMILES:} String representation of molecular structure
    \item \textbf{Graph:} Molecular graph with atoms as nodes and bonds as edges
    \item \textbf{3D coordinates:} 3D spatial coordinates of atoms
    \item \textbf{Descriptors:} Molecular descriptors and fingerprints
\end{itemize}

\subsection{Molecular Property Prediction}

Neural networks can predict molecular properties:
\begin{equation}
    p = f_\theta(\text{molecule})
\end{equation}

where $p$ is the property of interest (e.g., solubility, toxicity).

\subsection{Drug-Target Interaction}

Predict interactions between drugs and targets:
\begin{equation}
    \text{Interaction} = g_\theta(\text{drug}, \text{target})
\end{equation}

\subsection{Molecular Generation}

Generate new molecules with desired properties:
\begin{itemize}
    \item \textbf{VAE-based:} Use variational autoencoders
    \item \textbf{GAN-based:} Use generative adversarial networks
    \item \textbf{Flow-based:} Use normalizing flows
    \item \textbf{Diffusion-based:} Use diffusion models
\end{itemize}

\section{Materials Science}

\subsection{Crystal Structure Prediction}

Predict crystal structures from chemical composition:
\begin{equation}
    \text{Structure} = h_\theta(\text{composition})
\end{equation}

\subsection{Property Prediction}

Predict material properties:
\begin{itemize}
    \item \textbf{Electronic properties:} Band gap, conductivity
    \item \textbf{Mechanical properties:} Elastic modulus, strength
    \item \textbf{Thermal properties:} Thermal conductivity, heat capacity
    \item \textbf{Magnetic properties:} Magnetization, Curie temperature
\end{itemize}

\subsection{Materials Discovery}

Discover new materials with desired properties:
\begin{itemize}
    \item \textbf{High-throughput screening:} Screen large databases
    \item \textbf{Inverse design:} Design materials for specific properties
    \item \textbf{Multi-objective optimization:} Optimize multiple properties
\end{itemize}

\section{Climate Science}

\subsection{Weather Prediction}

Predict weather patterns using neural networks:
\begin{equation}
    \text{Weather}_{t+1} = f_\theta(\text{Weather}_t, \text{Forcing}_t)
\end{equation}

\subsection{Climate Modeling}

Model climate systems:
\begin{itemize}
    \item \textbf{Atmospheric dynamics:} Model atmospheric processes
    \item \textbf{Ocean dynamics:} Model ocean circulation
    \item \textbf{Carbon cycle:} Model carbon fluxes
    \item \textbf{Climate change:} Predict climate change impacts
\end{itemize}

\subsection{Extreme Weather Events}

Predict and analyze extreme weather events:
\begin{itemize}
    \item \textbf{Hurricanes:} Predict hurricane formation and intensity
    \item \textbf{Heat waves:} Predict heat wave occurrence
    \item \textbf{Droughts:} Predict drought conditions
    \item \textbf{Floods:} Predict flood risk
\end{itemize}

\section{Physics}

\subsection{Particle Physics}

AI applications in particle physics:
\begin{itemize}
    \item \textbf{Event classification:} Classify particle physics events
    \item \textbf{Track reconstruction:} Reconstruct particle tracks
    \item \textbf{Jet tagging:} Identify different types of jets
    \item \textbf{Anomaly detection:} Detect new physics signals
\end{itemize}

\subsection{Plasma Physics}

AI applications in plasma physics:
\begin{itemize}
    \item \textbf{Turbulence modeling:} Model plasma turbulence
    \item \textbf{Confinement prediction:} Predict plasma confinement
    \item \textbf{Control:} Control plasma parameters
    \item \textbf{Diagnostics:} Analyze plasma diagnostics
\end{itemize}

\subsection{Quantum Physics}

AI applications in quantum physics:
\begin{itemize}
    \item \textbf{Quantum state tomography:} Reconstruct quantum states
    \item \textbf{Quantum error correction:} Correct quantum errors
    \item \textbf{Quantum optimization:} Solve optimization problems
    \item \textbf{Quantum machine learning:} Use quantum computers for ML
\end{itemize}

\section{Biology}

\subsection{Protein Structure Prediction}

Predict protein structures from sequences:
\begin{equation}
    \text{Structure} = f_\theta(\text{Sequence})
\end{equation}

\subsection{Protein Function Prediction}

Predict protein functions:
\begin{itemize}
    \item \textbf{Enzyme classification:} Classify enzymes
    \item \textbf{Binding sites:} Predict binding sites
    \item \textbf{Protein-protein interactions:} Predict interactions
    \item \textbf{Drug targets:} Identify drug targets
\end{itemize}

\subsection{Genomics}

AI applications in genomics:
\begin{itemize}
    \item \textbf{Sequence analysis:} Analyze DNA/RNA sequences
    \item \textbf{Variant calling:} Identify genetic variants
    \item \textbf{Gene expression:} Predict gene expression
    \item \textbf{Epigenetics:} Analyze epigenetic modifications
\end{itemize}

\section{Scientific Computing}

\subsection{Numerical Methods}

AI can accelerate numerical methods:
\begin{itemize}
    \item \textbf{Finite element methods:} Accelerate FEM simulations
    \item \textbf{Monte Carlo methods:} Improve MC sampling
    \item \textbf{Optimization:} Solve optimization problems
    \item \textbf{Integration:} Numerical integration
\end{itemize}

\subsection{Surrogate Models}

AI can create surrogate models for expensive simulations:
\begin{equation}
    \text{Output} = f_\theta(\text{Input})
\end{equation}

where $f_\theta$ is much faster than the original simulation.

\subsection{Uncertainty Quantification}

Quantify uncertainty in AI predictions:
\begin{itemize}
    \item \textbf{Bayesian neural networks:} Bayesian uncertainty
    \item \textbf{Ensemble methods:} Ensemble uncertainty
    \item \textbf{Conformal prediction:} Distribution-free uncertainty
    \item \textbf{Calibration:} Calibrate uncertainty estimates
\end{itemize}

\section{Scientific Data Analysis}

\subsection{Data Preprocessing}

Preprocess scientific data:
\begin{itemize}
    \item \textbf{Noise reduction:} Remove noise from data
    \item \textbf{Outlier detection:} Detect and handle outliers
    \item \textbf{Data augmentation:} Augment limited datasets
    \item \textbf{Normalization:} Normalize data for analysis
\end{itemize}

\subsection{Feature Engineering}

Engineer features for scientific data:
\begin{itemize}
    \item \textbf{Domain knowledge:} Use domain-specific features
    \item \textbf{Physical constraints:} Enforce physical constraints
    \item \textbf{Symmetries:} Respect symmetries
    \item \textbf{Invariants:} Use invariant features
\end{itemize}

\subsection{Interpretability}

Make AI models interpretable for scientists:
\begin{itemize}
    \item \textbf{Feature importance:} Identify important features
    \item \textbf{Attention mechanisms:} Focus on relevant parts
    \item \textbf{Saliency maps:} Visualize model attention
    \item \textbf{Counterfactuals:} Generate counterfactual explanations
\end{itemize}

\section{Challenges and Limitations}

\subsection{Data Challenges}

\begin{itemize}
    \item \textbf{Limited data:} Scientific datasets are often small
    \item \textbf{Noisy data:} Experimental data can be noisy
    \item \textbf{Imbalanced data:} Classes may be imbalanced
    \item \textbf{Missing data:} Some data may be missing
\end{itemize}

\subsection{Model Challenges}

\begin{itemize}
    \item \textbf{Generalization:} Models may not generalize well
    \item \textbf{Interpretability:} Models may be hard to interpret
    \item \textbf{Uncertainty:} Quantifying uncertainty is challenging
    \item \textbf{Validation:} Validating models is difficult
\end{itemize}

\subsection{Integration Challenges}

\begin{itemize}
    \item \textbf{Workflow integration:} Integrating with existing workflows
    \item \textbf{Expert knowledge:} Incorporating expert knowledge
    \item \textbf{Validation:} Validating AI predictions
    \item \textbf{Trust:} Building trust in AI systems
\end{itemize}

\section{Recent Advances}

\subsection{Large Language Models for Science}

Large language models can help with scientific tasks:
\begin{itemize}
    \item \textbf{Literature review:} Summarize scientific literature
    \item \textbf{Hypothesis generation:} Generate scientific hypotheses
    \item \textbf{Code generation:} Generate scientific code
    \item \textbf{Data analysis:} Analyze scientific data
\end{itemize}

\subsection{Multimodal AI}

Multimodal AI can handle multiple data types:
\begin{itemize}
    \item \textbf{Text and images:} Combine text and image data
    \item \textbf{Sequences and structures:} Combine sequences and structures
    \item \textbf{Time series and images:} Combine time series and images
    \item \textbf{Multiple modalities:} Handle multiple data modalities
\end{itemize}

\subsection{Federated Learning}

Federated learning enables collaborative research:
\begin{itemize}
    \item \textbf{Privacy:} Protect sensitive data
    \item \textbf{Collaboration:} Enable collaboration across institutions
    \item \textbf{Scalability:} Scale to large datasets
    \item \textbf{Efficiency:} Reduce communication costs
\end{itemize}

\section{Applications}

\subsection{COVID-19 Research}

AI has been used extensively in COVID-19 research:
\begin{itemize}
    \item \textbf{Drug repurposing:} Identify existing drugs for COVID-19
    \item \textbf{Vaccine design:} Design vaccines
    \item \textbf{Epidemiological modeling:} Model disease spread
    \item \textbf{Medical imaging:} Analyze medical images
\end{itemize}

\subsection{Cancer Research}

AI applications in cancer research:
\begin{itemize}
    \item \textbf{Drug discovery:} Discover new cancer drugs
    \item \textbf{Personalized medicine:} Tailor treatments to patients
    \item \textbf{Medical imaging:} Analyze medical images
    \item \textbf{Biomarker discovery:} Discover new biomarkers
\end{itemize}

\subsection{Renewable Energy}

AI applications in renewable energy:
\begin{itemize}
    \item \textbf{Solar energy:} Optimize solar panel placement
    \item \textbf{Wind energy:} Predict wind patterns
    \item \textbf{Energy storage:} Optimize energy storage
    \item \textbf{Grid management:} Manage power grids
\end{itemize}

\section{Future Directions}

\subsection{Theoretical Advances}

\begin{itemize}
    \item \textbf{Physics-informed neural networks:} Incorporate physics into neural networks
    \item \textbf{Uncertainty quantification:} Better uncertainty quantification
    \item \textbf{Interpretability:} More interpretable models
    \item \textbf{Generalization:} Better generalization across domains
\end{itemize}

\subsection{Methodological Improvements}

\begin{itemize}
    \item \textbf{Better architectures:} Design better architectures for science
    \item \textbf{Training techniques:} Improve training techniques
    \item \textbf{Data augmentation:} Better data augmentation methods
    \item \textbf{Transfer learning:} Better transfer learning across domains
\end{itemize}

\subsection{Applications}

\begin{itemize}
    \item \textbf{New domains:} Apply to new scientific domains
    \item \textbf{Real-world deployment:} Deploy in real-world applications
    \item \textbf{Collaboration:} Enable better collaboration between AI and scientists
    \item \textbf{Education:} Use AI for scientific education
\end{itemize}

\section{Conclusion}

AI for science represents a transformative paradigm in scientific research:

\textbf{Key contributions:}
\begin{itemize}
    \item \textbf{Accelerated discovery:} Accelerate scientific discovery
    \item \textbf{New insights:} Generate new scientific insights
    \item \textbf{Efficiency:} Improve research efficiency
    \item \textbf{Collaboration:} Enable new forms of collaboration
\end{itemize}

\textbf{Key insights:}
\begin{itemize}
    \item \textbf{Data integration:} Integrate diverse scientific data
    \item \textbf{Physical constraints:} Respect physical constraints
    \item \textbf{Uncertainty:} Quantify and propagate uncertainty
    \item \textbf{Interpretability:} Make models interpretable for scientists
\end{itemize}

\textbf{Future work:}
\begin{itemize}
    \item \textbf{Theoretical understanding:} Better theoretical understanding
    \item \textbf{Methodological advances:} New methodological advances
    \item \textbf{Applications:} New applications and domains
    \item \textbf{Collaboration:} Better collaboration between AI and scientists
\end{itemize}

AI for science continues to drive innovation in scientific research, with ongoing work exploring new applications, methodologies, and theoretical foundations.


% ==== Back Matter ====
\backmatter

% Bibliography
\bibliographystyle{plain}
\bibliography{references}

% Index
\printindex

\end{document}
